\documentclass[11pt]{article}
\usepackage[T1]{fontenc}
\usepackage{lmodern}
\usepackage[letterpaper,margin=0.88in]{geometry}
\usepackage{amsmath,amssymb,amsthm,mathtools}
\usepackage{booktabs,longtable,array,enumitem,microtype,xcolor,xurl}
\usepackage{fancyhdr}
\usepackage[colorlinks=true,linkcolor=navy,citecolor=navy,urlcolor=navy]{hyperref}
\definecolor{navy}{HTML}{163B5C}
\definecolor{muted}{HTML}{536575}
\definecolor{accent}{HTML}{8A521A}
\newtheorem{theorem}{Theorem}[section]
\newtheorem{lemma}[theorem]{Lemma}
\newtheorem{proposition}[theorem]{Proposition}
\newtheorem{corollary}[theorem]{Corollary}
\theoremstyle{definition}
\newtheorem{definition}[theorem]{Definition}
\newtheorem{example}[theorem]{Example}
\newtheorem{conjecture}[theorem]{Conjecture}
\theoremstyle{remark}
\newtheorem{remark}[theorem]{Remark}
\numberwithin{equation}{section}
\DeclareMathOperator{\Li}{Li}
\DeclareMathOperator{\FP}{FP}
\DeclareMathOperator{\CT}{CT}
\DeclareMathOperator{\Res}{Res}
\newcommand{\dd}{\,\mathrm d}
\newcommand{\ii}{\mathrm i}
\newcommand{\Q}{\mathbb Q}
\newcommand{\C}{\mathbb C}
\newcommand{\N}{\mathbb N}
\newcommand{\re}{\operatorname{Re}}
\newcommand{\im}{\operatorname{Im}}
\newcommand{\src}[1]{\texttt{\small\detokenize{#1}}}
\setlist{leftmargin=*,itemsep=3pt,topsep=5pt}
\setlength{\emergencystretch}{2em}
\setlength{\headheight}{14pt}
\pagestyle{fancy}
\fancyhf{}
\fancyhead[L]{\small\color{muted}Exact harmonic-tail identities}
\fancyhead[R]{\small\color{muted}ProveIt research continuation}
\fancyfoot[C]{\thepage}
\renewcommand{\headrulewidth}{0.3pt}
\hypersetup{pdftitle={Exact Harmonic-Tail Identities: Gap-Four Reconstruction, Hurwitz Shifts, and Alternating Moments},pdfauthor={Research continuation prepared for Vladimir Reshetnikov},pdfsubject={Exact identities for harmonic tails, polygamma, polylogarithms and Stieltjes constants}}

\begin{document}
\hypersetup{pageanchor=false}
\begin{titlepage}
\centering
\vspace*{0.9in}
{\Large\color{navy}\sffamily PROVEIT RESEARCH CONTINUATION\par}
\vspace{0.45in}
{\LARGE\bfseries Exact Harmonic-Tail Identities\par}
\vspace{0.2in}
{\Large Gap-Four Reconstruction, Hurwitz Shifts,\\[5pt]
and Alternating Moments\par}
\vspace{0.5in}
{\large Prepared for Vladimir Reshetnikov\par}
\vspace{0.15in}
{11 October 2026\par}
\vspace{0.55in}
\begin{minipage}{0.89\textwidth}
\small
\textbf{Abstract.}
We continue the exact harmonic-tail programme in the ProveIt
polylogarithm manuscripts. For every pair of even orders with
$q-p=4$, the first $q$ centered moments generate the full specified
formal coordinate space. The proof gives an explicit nonzero
determinant and a reconstruction requiring only a two-by-two inverse.
Its central comparison theorem proves that every generated minor of
order two in the fundamental recurrence is strictly negative for
arbitrary even gaps. A concrete new $(8,12)$ identity is also certified
by a finite rational telescoper.

We then develop exact shifted Hurwitz-tail moments, including the
Bernoulli corrections that distinguish ordinary subtracted sums from
continued Dirichlet values. An elementary trigamma-square telescoper
gives a free-shift identity and an all-order differentiated family in
ordinary zeta values. Spectral coefficients retain the required
generalized Stieltjes terms. Finally, signed tail products lead to
alternating moment reductions, a two-parameter digamma generating
identity, and anchored polylogarithmic primitives. All analytic
statements carry explicit convergence or continuation conventions.
The accompanying programs check finite rational identities and
independent numerical representations; the infinite-family claims
are proved in the text.
\end{minipage}
\vfill
\begin{minipage}{0.89\textwidth}
\footnotesize\color{muted}
Repository baseline:
\texttt{446dca9bd65c20e409d1d2561d4ffad408103e2a}.\\
The manuscript source files and all six incoming archives were fixed
to this revision. The new conclusions are separated from inherited
results and from questions that remain open.
\end{minipage}
\end{titlepage}
\clearpage
\pagenumbering{roman}
\hypersetup{pageanchor=true}
\tableofcontents
\clearpage
\pagenumbering{arabic}
\section{The questions advanced in this continuation}
\label{sec:scope}

The harmonic tails
\[
 A_p(n)=\zeta(p)-H_n^{(p)}=\zeta(p,n+1),\qquad
 H_n^{(p)}=\sum_{k=1}^n k^{-p},\quad H_0^{(p)}=0,
\]
connect three languages used throughout the existing programme:
ordinary sums involving harmonic numbers, polygamma functions at
integer or shifted arguments, and multiple zeta or polylogarithm
values. A polynomial weight can destroy ordinary convergence without
destroying the exact identity. The useful question is then to specify
the finite subtraction and its constant, and to determine which
finite collection of moments already contains all the information.

The present article concentrates on this exact-identity question.
Asymptotic expansions enter only to justify convergence, constant
extraction, and limiting operations. The main result settles one
additional infinite family of the existing optimal-prefix conjecture.
The other results produce explicit shift, twist, derivative, and
primitive formulas around the same tail functions.

\subsection{The precise change in the conjecture}

For even $2\le p\le q$ consider the continued centered moments
\[
 D_{p,q}(-2r),\qquad
 D_{p,q}(s)=\sum_{n\ge0}
       A_p(n)A_q(n)(n+\tfrac12)^{-s}.
\]
The source reports supply $q$ formal coordinates and prove that
some finite initial segment generates them. Their stronger conjecture
asks whether the first $q$ moments always suffice. At the pinned
revision the following families have proofs in the source reports:
the diagonal $p=q$, the adjacent case $q-p=2$, and each fixed smaller
order $p=2,4,6$ \cite{Canonical,Differential,Bernoulli,Polygamma}.

Theorem~\ref{eg:main} proves the entire family $q-p=4$.
Its additional cases begin at $(p,q)=(8,12)$, since the first three
pairs overlap the fixed-order results. The proof is an all-order
nonvanishing argument. A table of computed determinants is only a
normalization check. Theorem~\ref{eg:two-minors} supplies the underlying
comparison theorem for every even gap, while the text explains why
that theorem alone does not settle higher-dimensional determinants.

The rank is taken over $\Q$ on the source's specified formal labels.
It yields exact reconstruction identities for the corresponding
numbers, whether or not those numbers have further arithmetic
relations. It does not assert linear independence of zeta values.

\begin{center}
\begin{tabular}{@{}>{\raggedright\arraybackslash}p{0.25\textwidth}>{\raggedright\arraybackslash}p{0.31\textwidth}>{\raggedright\arraybackslash}p{0.36\textwidth}@{}}
\toprule
Direction & Earlier status & Contribution here \\
\midrule
Optimal moment prefix & Diagonal, gap two, and $p=2,4,6$ proved &
Gap four proved at every order; explicit reconstruction and sums \\
Hurwitz shift & Requested extension of half-integer centering &
Finite endpoint and Bernoulli corrections; free-shift trigamma identity \\
Logarithmic moments & Subtraction constants must be tracked &
All spectral coefficients, including Stieltjes resonance terms \\
Character twists & Proposed exact reduction programme &
Signed centered moments, digamma generator, polylogarithmic primitives \\
Gaussian $S_6$, revised $S_8$ & Numerical candidates with unresolved proofs &
Remain open; no claim of resolution \\
\bottomrule
\end{tabular}
\end{center}

\subsection{Source scope and established mathematics}

The repository snapshot is \cite{Repo}. The current incoming directory
contains six reports: \emph{Bernoulli--Tornheim--Digamma},
\emph{Cubic Tail Descent}, \emph{Determinantal Harmonic Resonance},
\emph{Differential Polygamma and Lerch Jets},
\emph{Multishift Reflection}, and \emph{Polygamma Tail Identities}.
The source manifest records the Git blob of every acquired manuscript
text and archive. We inspected the claims used here and their nearby
research agendas; this is a targeted proof audit, not a certification
of every theorem in the collective volumes.

Several ingredients are established mathematics. Hoffman's treatment
of products of Riemann-zeta tails explains their relation to multiple
zeta values and higher products \cite{Hoffman}. Teo proves the
odd-weight reduction of alternating double Euler sums \cite{Teo}.
The Hurwitz recurrence, its relation to polygamma, the Bernoulli
difference law, and Euler--Maclaurin continuation are classical
\cite{DLMF}. Bernoulli Hankel determinants also have an extensive
literature \cite{DilcherJiu}. These ingredients are credited where
used. The additional contribution is the specified determinant proof,
the normalized extensions, and their explicit consequences relative
to the inspected programme. An exhaustive historical-priority claim
is not made.

\subsection{Conventions needed for reproducibility}

Bernoulli numbers use $B_1=-1/2$, and Bernoulli polynomials satisfy
\[
 \frac{te^{xt}}{e^t-1}=\sum_{m\ge0}B_m(x)\frac{t^m}{m!}.
\]
The multiple-polylogarithm convention uses decreasing indices:
\[
 \Li_{s_1,\ldots,s_d}(z_1,\ldots,z_d)
 =\sum_{n_1>\cdots>n_d\ge1}
       \frac{z_1^{n_1}\cdots z_d^{n_d}}
            {n_1^{s_1}\cdots n_d^{s_d}}.
\]
In particular, a one-variable depth-two expression
$\Li_{p,q}(z,1)$ weights the larger index. The generalized Stieltjes
constants use
\[
 \zeta(1+s,a)=\frac1s+
       \sum_{j\ge0}\frac{(-1)^j\gamma_j(a)}{j!}s^j,
 \qquad \gamma_0(a)=-\psi(a).
\]
A superscript $\zeta^{(j)}$ denotes a derivative in the spectral
variable. Derivatives of $\psi$ denote derivatives in its argument.
Every finite part and every convergent subtraction is specified at
its point of use. No divergent series is assigned a value without
a stated continuation or subtraction rule.

\section{Centered moments and their exact normalization}
\label{sec:foundation}

We first establish the analytic meaning of the moment values used in
the reconstruction theorem. The finite coordinate formulas are
inherited from the tail reports; their normalization is reproduced
here so that the additional determinant argument can be checked
without reconstructing an earlier archive.

\subsection{Centered asymptotic subtraction}

Put $x_n=n+1/2$. For an integer $p\ge2$, the centered Hurwitz
expansion is
\begin{equation}\label{found:asymptotic}
 A_p(n)\sim\sum_{j\ge0}a_j(p)x_n^{1-p-2j},\qquad
 a_0(p)=\frac1{p-1},\quad
 a_j(p)=\frac{B_{2j}(1/2)(p)_{2j-1}}{(2j)!}\quad(j\ge1).
\end{equation}
Here $(p)_m=p(p+1)\cdots(p+m-1)$ is the rising factorial.
For each fixed truncation the usual Euler--Maclaurin remainder
justifies this expansion, locally uniformly in any parameters kept
away from their poles \cite{DLMF}. In particular, if $p,q$ are even,
\begin{equation}\label{found:product-expansion}
 A_p(n)A_q(n)\sim
 \sum_{j\ge0}c_j(p,q)x_n^{2-p-q-2j},\qquad
 c_j(p,q)=\sum_{h=0}^j a_h(p)a_{j-h}(q).
\end{equation}

\begin{proposition}[Continued and ordinary centered moments]
\label{found:regularity}
For even $p,q\ge2$, $D_{p,q}(s)$ initially converges absolutely for
$\Re s>3-p-q$ and has a meromorphic continuation to the plane.
Its possible poles are among
\[
                         s=3-p-q-2j,\qquad j\ge0.
\]
Every $s=-2r$, $r\ge0$, is regular. Moreover, with
$J=r+1-(p+q)/2$, the exact ordinary identity is
\begin{equation}\label{found:ordinary}
 D_{p,q}(-2r)=\sum_{n\ge0}\left\{
 x_n^{2r}A_p(n)A_q(n)
 -\sum_{j=0}^{J}c_j(p,q)x_n^{2r+2-p-q-2j}\right\}.
\end{equation}
The inner sum is empty when $J<0$. The complete summand is
$O(x_n^{-2})$ when subtraction is needed; when it is not needed the
original series is absolutely convergent.
\end{proposition}

\begin{proof}
Subtract $L$ terms of \eqref{found:product-expansion}. Summing them
back gives
\begin{align}\label{found:continuation}
 D_{p,q}(s)={}&\sum_{n\ge0}
 \left(A_p(n)A_q(n)-\sum_{j=0}^{L-1}c_jx_n^{2-p-q-2j}\right)x_n^{-s}
 \notag\\
 &+\sum_{j=0}^{L-1}c_j\zeta(s+p+q-2+2j,1/2).
\end{align}
The first term converges normally for
$\Re s>3-p-q-2L$. Increasing $L$ gives the continuation and the
listed possible poles. Their positions are odd integers, so none
equals $-2r$. At this value choose exactly the nonnegative powers
in the polynomial subtraction. Each added-back Hurwitz value is
zero because
\[
 \zeta(-2h,1/2)=-\frac{B_{2h+1}(1/2)}{2h+1}=0
 \qquad(h\ge0).
\]
This proves \eqref{found:ordinary}. The statement uses a convergent
sum of the complete braces; the separately divergent pieces are
never summed on their own.
\end{proof}

The last vanishing is special to the chosen center. It also cannot
be differentiated as though it were an identity in the spectral
variable: for example $\zeta'(0,1/2)=-\tfrac12\log2\ne0$.
The shifted and logarithmic formulas below preserve the corresponding
finite correction terms.

\subsection{Finite summation and the block transform}

Let $\CT_N$ denote the coefficient of $N^0(\log N)^0$ in a
specified large-$N$ expansion. Centered even power sums satisfy
\begin{equation}\label{found:power-sum}
 S_r(N)=\sum_{n=0}^{N-1}(n+1/2)^{2r}
 =\frac{B_{2r+1}(N+1/2)}{2r+1}
 =\sum_{j=0}^r T_{rj}N^{2j+1},
\end{equation}
where
\[
 T_{rj}=\frac{\binom{2r}{2j}}{2j+1}B_{2r-2j}(1/2).
\]
The polynomial has zero constant term. By
Proposition~\ref{found:regularity}, the moment is consequently the
constant term of its finite partial sum. The recurrence
\[
                         A_p(k-1)=A_p(k)+k^{-p}
\]
and finite summation by parts give
\begin{equation}\label{found:blocks}
 D_{p,q}(-2r)=\sum_{j=0}^r T_{rj}C_{p,q}(2j+1),
\end{equation}
where, for a positive odd $d$,
\begin{align}\label{found:block-CT}
 C_{p,q}(d)=\CT_N\bigg\{&N^dA_p(N)A_q(N)
       +\sum_{k=1}^N k^{d-q}A_p(k)\notag\\
       &+\sum_{k=1}^N k^{d-p}A_q(k)
       +\sum_{k=1}^N k^{d-p-q}\bigg\}.
\end{align}
This identity is obtained at a finite cutoff before any limiting
operation. It is also the convenient starting point for shifted and
twisted extensions.

For later reference we state the explicit inherited evaluation.
Set $W=p+q$, and define
\begin{align}
 U_{p,q}(d)&=
 \begin{cases}
 \zeta(p,q-d)+\tfrac12\zeta(W-d),&d<q,\\[3pt]
 \displaystyle\sum_{h=1}^{\min\{(d-q+1)/2,p/2-1\}}
 \frac{\binom{d-q+1}{2h}}{d-q+1}
 B_{d-q+1-2h}\zeta(p-2h),&d>q,
 \end{cases}\label{found:U}\\
 R_{p,q}(d)&=
 \begin{cases}
 0,&d<W-1,\\[3pt]
 \displaystyle-\frac{B_{d-W+1}}2
 \left\{\frac1p\binom{d-q}{p-1}
       +\frac1q\binom{d-p}{q-1}\right\},&d\ge W-1.
 \end{cases}\label{found:R}
\end{align}
Then
\begin{equation}\label{found:C-evaluation}
 C_{p,q}(d)=U_{p,q}(d)+U_{q,p}(d)+R_{p,q}(d).
\end{equation}
All double-zeta sums appearing in \eqref{found:U} are convergent.
There is no case $d=q$ because $d$ is odd and $q$ is even.

\paragraph{Why these finite formulas hold.}
For $d<q$, changing the order in the convergent sum yields
$\sum_{k\ge1}k^{d-q}A_p(k)=\zeta(p,q-d)$.
For $d>q$ set $m=d-q$ and
$F_m(N)=(B_{m+1}(N+1)-B_{m+1})/(m+1)$. The finite identity
\[
 \sum_{k=1}^N k^mA_p(k)
 =A_p(N)F_m(N)+
                  \sum_{\ell=1}^N\frac{F_m(\ell-1)}{\ell^p}
\]
reduces everything to Bernoulli polynomials and ordinary power sums.
The even Bernoulli terms give \eqref{found:U}. The half-leading Faulhaber terms leave precisely
one half of $\zeta(W-d)$ for each convergent ordered block. When
both ordered blocks require reduction, their two half-terms cancel
the last sum in \eqref{found:block-CT}. At $d=W-1$ this includes cancellation of
the two $-\gamma/2$ constants against $+\gamma$ from the harmonic
sum. The remaining boundary products give \eqref{found:R}.
This is the finite-sum proof in the incoming tail calculus
\cite{Canonical,Bernoulli}; no additional relation among numerical
periods is used.

\subsection{Inverse transform and a small check}

Write $C_j=C_{p,q}(2j+1)$ and $D_r=D_{p,q}(-2r)$.
The Bernoulli generating identity in \eqref{found:power-sum} yields
the inverse transform
\begin{equation}\label{found:inverse}
 C_j=\sum_{r=0}^j\binom{2j+1}{2r}2^{-2(j-r)}D_r.
\end{equation}
One way to verify it is to set
$G(t)=\sum D_rt^{2r}/(2r)!$ and
$C(t)=\sum C_jt^{2j+1}/(2j+1)!$. The forward transform says
$G(t)=C(t)/(2\sinh(t/2))$ as a formal power series. Multiplying
back proves \eqref{found:inverse}, with no analytic assumption on
the growth of the moments.

For $p=q=2$, the first two blocks are $C_0=3\zeta(3)$ and
$C_1=-1/2$. Hence
\[
 D_0=3\zeta(3),\qquad
 D_1=-\frac14\zeta(3)-\frac16,
\]
and Proposition~\ref{found:regularity} gives the ordinary identity
\[
 \sum_{n\ge0}\bigl((n+1/2)^2\psi'(n+1)^2-1\bigr)
 =-\frac14\zeta(3)-\frac16.
\]
These are inherited checks. The next section uses the same fixed
normalization to prove an additional infinite reconstruction family.

% Integration fragment. All labels start eg: and no special macros are required.
% Prior source: Differential_Polygamma_and_Lerch_Identities, Section 8;
% current Bernoulli_Tornheim_Digamma_Identities, Sections 3--6.
% Exact verification: verify_gap_four.py, gap_four_exact_results.json.
\section{The optimal initial prefix when the even orders differ by four}
\label{eg:section}

The current incoming reports prove the optimal initial-prefix statement
for the diagonal, for adjacent even orders, and for the fixed smaller
orders $p=2,4,6$ \cite{Canonical,Differential,Bernoulli,Polygamma}.
We prove the next fixed-gap family
\[
                    (p,q)=(2b-4,2b),\qquad b\ge3.
\]
Its additional range relative to those reports begins with $(p,q)=(8,12)$ and continues
with $(10,14),(12,16),\ldots$. The proof uses the strict log-convexity
of the positive coefficients of $2/x-\cot(x/2)$ to show that a
specified two-by-two boundary determinant never vanishes. We also
prove a structural statement about every pair of columns in the
fundamental recurrence for an arbitrary even gap. The latter does
not settle the higher-dimensional boundary determinant.

\subsection{Normalization and the determinant to be proved nonzero}

For even $2\le p\le q=2b$, put
\[
 A_p(n)=\zeta(p)-H_n^{(p)}=\zeta(p,n+1),\qquad
 D_{p,q}(s)=\sum_{n\ge0}A_p(n)A_q(n)(n+\tfrac12)^{-s}.
\]
The series initially converges for $\Re s>3-p-q$, and its continuation
is regular at every $s=-2r$. Write $D_r=D_{p,q}(-2r)$. We use the
inherited blocks $C_j$, characterized by the mutually inverse formulas
\begin{align}
 D_r&=\sum_{j=0}^r\frac{\binom{2r}{2j}}{2j+1}
       B_{2r-2j}(\tfrac12)C_j,\label{eg:forward}\\
 C_j&=\sum_{r=0}^j\binom{2j+1}{2r}2^{-2(j-r)}D_r.
       \label{eg:backward}
\end{align}
For $j\ge b$ they have the finite coordinate form
\[
 C_j=\sum_{k=0}^{b-1}a_{jk}Z_k,\qquad
 Z_k=\zeta(2k),\qquad Z_0=-\tfrac12.
\]
For $d=2j+1$ and $e=d-p-q+2k+1$, the inherited coefficients are
\begin{equation}\label{eg:entries}
 a_{jk}=\begin{cases}
 B_e\left\{\displaystyle
 \frac{\binom{d-q+1}{p-2k}}{d-q+1}
 +\frac{\binom{d-p+1}{q-2k}}{d-p+1}\right\},&e\ge0,\\[6pt]
 0,&e<0.
 \end{cases}
\end{equation}
Each binomial summand is omitted unless its lower index lies between
$2$ and its upper index. In particular a zero lower index is omitted.
Let
\begin{equation}\label{eg:M}
 M_{p,q}=(a_{b+r-1,k})_{1\le r\le b,\ 0\le k<b}.
\end{equation}
The first $q$ moments have full rank on the specified formal labels
\[
 C_0,\ldots,C_{b-1},Z_0,\ldots,Z_{b-1}
\]
if and only if $M_{p,q}$ is nonsingular. The complete first-$q$
determinant in this coordinate order is
$\det M_{p,q}/(4b-1)!!$, by the triangular transform
\eqref{eg:forward}. No arithmetic independence of the actual constants
is asserted by this formal rank statement.

\subsection{A strict comparison theorem for all even gaps}

Put $q-p=2c$, where $1\le c<b$, and define
\begin{equation}\label{eg:alpha}
 \alpha_n=\frac{2|B_{2n}|}{(2n)!}
          =4\sum_{m=1}^{\infty}(2\pi m)^{-2n},\qquad n\ge1.
\end{equation}
For $s\ge c+1$ set
\begin{equation}\label{eg:ABD}
 A_s=\frac{(2s)!}{(2s-2c)!},\qquad
 B_s=\frac{(2s-1)!}{(2s-2c-1)!},\qquad
 d_s=2(A_s+B_s),
\end{equation}
and put $A_s=0$ for $1\le s\le c$. The letters $A_s,B_s$ in
\eqref{eg:ABD} are rational recurrence coefficients, distinguished from
the harmonic tails $A_p(n)$ and Bernoulli numbers $B_{2n}$ by their
arguments and subscripts. Define
\begin{equation}\label{eg:kernel}
 K(s,u)=\frac{\alpha_{s-u}(B_s+A_u)}{d_s},
 \qquad s\ge c+1,\quad1\le u<s,
\end{equation}
and the fundamental array
\begin{align}
 R_{s,j}&=\delta_{sj},&&1\le s,j\le c,\notag\\
 R_{s,j}&=\sum_{u=1}^{s-1}K(s,u)R_{u,j},
       &&s\ge c+1,\quad1\le j\le c.
       \label{eg:fundamental}
\end{align}
All generated entries $R_{s,j}$, $s\ge c+1$, are strictly positive.
The earlier fixed-gap reduction gives
\begin{equation}\label{eg:inherited-reduction}
 \det M_{2(b-c),2b}
 =\frac{(-1)^{b(b-1)/2}}{4^bb!}
    \left(\prod_{s=c+1}^{b+c}d_s\right)
    \det(R_{b+i,j})_{1\le i,j\le c}.
\end{equation}
We give the coefficient mechanism below to check the convention and
signs directly; nonvanishing of the last determinant is the additional
issue addressed here when $c=2$.

\begin{lemma}[Strictly decreasing kernel ratios]\label{eg:kernel-ratios}
For every $s\ge c+1$, the quotient
\[
                   \frac{K(s+1,u)}{K(s,u)}
\]
is a strictly decreasing function of the integer $u$, $1\le u<s$.
\end{lemma}

\begin{proof}
The coefficients in \eqref{eg:alpha} are strictly log-convex. Indeed,
with $x_m=(2\pi m)^{-2}$,
\[
 \alpha_n\alpha_{n+2}-\alpha_{n+1}^2
 =16\sum_{m<\ell}x_m^nx_\ell^n(x_m-x_\ell)^2>0.
\]
The sums converge absolutely for $n\ge1$, and the positive support
contains more than one point. Thus $\alpha_{n+1}/\alpha_n$ strictly
increases with $n$. On the other hand,
\[
 \frac{K(s+1,u)}{K(s,u)}
 =\frac{d_s}{d_{s+1}}
   \frac{\alpha_{s+1-u}}{\alpha_{s-u}}
   \frac{B_{s+1}+A_u}{B_s+A_u}.
\]
The middle factor strictly decreases with $u$. The last factor is
nonincreasing because $A_u$ is nondecreasing, $B_{s+1}>B_s>0$,
and $(B_{s+1}+a)/(B_s+a)$ decreases as $a$ increases. Every factor
is positive, so their product strictly decreases.
\end{proof}

\begin{theorem}[Every generated minor of order two is negative]
\label{eg:two-minors}
Fix an even gap $2c$ with $c\ge2$. For any initial columns
$1\le j<k\le c$, the positive ratios
\[
                   \frac{R_{s,k}}{R_{s,j}},\qquad s\ge c+1,
\]
strictly decrease with $s$. Equivalently, for $c+1\le r<s$,
\begin{equation}\label{eg:minor-sign}
 \det\begin{pmatrix}R_{r,j}&R_{r,k}\\R_{s,j}&R_{s,k}\end{pmatrix}<0.
\end{equation}
\end{theorem}

\begin{proof}
Write $X_u=R_{u,j}$, $Y_u=R_{u,k}$. Suppose that the stated ratio
ordering has been proved through row $s\ge c+1$, and put
$z_s=Y_s/X_s$. This supposition is empty at $s=c+1$, so it includes
the induction base. For $u<s$ set $v_u=Y_u-z_sX_u$. Among the initial
rows, $v_j=-z_s<0$, $v_k=1>0$, and all other $v_u$ vanish. Among
the generated rows $c+1\le u<s$, the induction hypothesis gives
$v_u>0$. Thus the unique negative term is at index $j$, and every
positive term has larger index. The recurrence gives
\[
                     \sum_{u=1}^{s-1}K(s,u)v_u=0.
\]
Put $\beta_u=K(s+1,u)/K(s,u)$. Since $v_s=0$,
\begin{align*}
 Y_{s+1}-z_sX_{s+1}
 &=\sum_{u=1}^{s-1}K(s,u)\beta_uv_u\\
 &=\sum_{u=1}^{s-1}K(s,u)(\beta_u-\beta_j)v_u<0.
\end{align*}
The negative-index term vanishes; every nonzero remaining term is
negative by Lemma~\ref{eg:kernel-ratios}. Strictness already follows
from the term $u=k$. Hence $Y_{s+1}/X_{s+1}<z_s$, completing the
induction. Multiplying the ratio inequality by the two positive
denominators proves \eqref{eg:minor-sign}.
\end{proof}

The one exceptional initial negative index is essential in this proof.
Positivity of the recurrence entries alone would not imply the theorem.

\subsection{Exact determinant and reconstruction for gap four}

Now take $c=2$. Define $X_s=R_{s,1}$ and $Y_s=R_{s,2}$, so
\[
 (X_1,Y_1)=(1,0),\qquad (X_2,Y_2)=(0,1),
\]
and compute both columns by \eqref{eg:fundamental}. In this case
\begin{equation}\label{eg:d-gapfour}
 d_s=16(s-1)^2(2s-1)(2s-3),\qquad s\ge3.
\end{equation}
Set
\begin{equation}\label{eg:casoratian}
 E_s=X_{s+1}Y_s-X_sY_{s+1},\qquad s\ge3.
\end{equation}
Theorem~\ref{eg:two-minors} proves $E_s>0$ at every index. For example,
\[
 E_3=\frac1{14515200},\qquad
 E_4=\frac{61}{1170505728000},\qquad
 E_5=\frac{3419}{69528040243200000}.
\]
These are exact rational recurrence values, not approximations.

\begin{theorem}[Optimal prefix for every gap-four pair]\label{eg:main}
For every integer $b\ge3$,
\begin{equation}\label{eg:determinant}
 \boxed{\displaystyle
 \det M_{2b-4,2b}
 =(-1)^{b(b-1)/2+1}
   \frac{4^b((b+1)!)^2(2b+3)!!(2b+1)!!}{3b!}\,E_{b+1}\ne0.}
\end{equation}
Consequently, for $(p,q)=(2b-4,2b)$,
\begin{align}\label{eg:span}
 \operatorname{span}_{\mathbb Q}\{D_0,D_1,\ldots\}
 &=\operatorname{span}_{\mathbb Q}\{D_0,\ldots,D_{q-1}\}\notag\\
 &=\operatorname{span}_{\mathbb Q}
     \{C_0,\ldots,C_{b-1},Z_0,\ldots,Z_{b-1}\}.
\end{align}
The complete first-$q$ coefficient determinant is
\eqref{eg:determinant} divided by $(4b-1)!!$. The prefix is optimal
for rank $q$ on these $q$ formal labels.
\end{theorem}

\begin{proof}
The last determinant in \eqref{eg:inherited-reduction} is now
\[
 \det\begin{pmatrix}X_{b+1}&Y_{b+1}\\X_{b+2}&Y_{b+2}\end{pmatrix}
 =-E_{b+1}<0.
\]
The elementary product evaluation
\[
 \prod_{s=3}^{b+2}d_s
 =\frac{16^b((b+1)!)^2(2b+3)!!(2b+1)!!}{3}
\]
gives \eqref{eg:determinant}. The block decomposition and invertible
triangular transforms give \eqref{eg:span}. Fewer than $q$ rows cannot
have rank $q$ on $q$ specified formal coordinates.
\end{proof}

The first five determinants are
\[
\begin{array}{c|c}
(p,q)&\det M_{p,q}\\\hline
(2,6)&61/5760\\
(4,8)&-3419/138240\\
(6,10)&-713039/6451200\\
(8,12)&263589503/309657600\\
(10,14)&818219064067/78033715200
\end{array}
\]
The first three pairs overlap the previously settled fixed smaller
orders. The last two belong to the additional infinite range covered
by Theorem~\ref{eg:main}.

\paragraph{Checking the source reduction and the rotation.}
For an arbitrary coefficient vector $v=(v_0,\ldots,v_{b-1})$, set
$Q(t)=\sum_{k=0}^{b-1}v_k t^{2b-2k}/(2b-2k)!$ and write
$Q(ix)=R(x)=\sum_{u=1}^b r_ux^{2u}$. Thus
$r_u=(-1)^uv_{b-u}/(2u)!$. Consider the real odd expression
\[
 (R^{(2c)}-R^{(2c)}(0))\cot(x/2)+(R\cot(x/2))^{(2c)}.
\]
The coefficient at $x^{2r-1}$, with $s=r+c$, is
\begin{equation}\label{eg:coefficient-equation}
 Y_r^*=d_sr_s-\sum_{u=1}^{s-1}\alpha_{s-u}(B_s+A_u)r_u,
 \qquad r_{b+1}=\cdots=r_{b+c}=0.
\end{equation}
Indeed $2/x$ supplies $d_s$, and the remaining cotangent coefficients
are $-\alpha_n$. The constant derivative subtraction sets $A_c=0$,
exactly enforcing the omitted zero lower binomial index.
Comparing factorials directly with \eqref{eg:entries} gives the row scaling
\[
 Y_r^*=\frac{2(-1)^{r-c}}{(2r-1)!}(M_{p,q}v)_r.
\]
Here is a complete determinant check of
\eqref{eg:inherited-reduction}. Let $L$ be the lower triangular
matrix of size $b+c$ whose first $c$ rows are the corresponding
identity rows, and whose remaining entries are
\[
 L_{ss}=d_s,\qquad L_{su}=-d_sK(s,u)\ (u<s),\qquad
 L_{su}=0\ (u>s).
\]
Its inverse has entries $(L^{-1})_{s,j}=R_{s,j}$ in its first $c$
columns. The coefficient matrix $H$ of
\eqref{eg:coefficient-equation} is the minor with rows
$c+1,\ldots,b+c$ and columns $1,\ldots,b$. Jacobi's complementary
minor identity therefore gives
\[
 \det H=(-1)^{bc}\left(\prod_{s=c+1}^{b+c}d_s\right)
                     \det(R_{b+i,j})_{1\le i,j\le c}.
\]
If $T$ denotes the reversed diagonal map from $v$ to $r$ and $S$
the row scaling just displayed, then $HT=SM_{p,q}$ and
\[
 \det T=\frac{(-1)^b}{\prod_{u=1}^b(2u)!},\qquad
 \det S=\frac{2^b(-1)^{b(b+1)/2-bc}}
                       {\prod_{r=1}^b(2r-1)!}.
\]
Taking determinants yields exactly
\eqref{eg:inherited-reduction}. This checks the recurrence against
the original Bernoulli matrix, including the sign convention, without
a new analytic assumption.

For $c=2$, these equations give a particularly small reconstruction.
Given the actual blocks $C_b,\ldots,C_{2b-1}$, set
\[
 Y_r^*=\frac{2(-1)^{r-2}}{(2r-1)!}C_{b+r-1},\qquad1\le r\le b.
\]
Set $\eta_1=\eta_2=0$ and compute, for $3\le s\le b+2$,
\begin{equation}\label{eg:inhomogeneous}
 \eta_s=\frac{Y_{s-2}^*}{d_s}
          +\sum_{u=1}^{s-1}K(s,u)\eta_u.
\end{equation}
The two initial coefficients are recovered by
\begin{equation}\label{eg:reconstruction}
 \begin{pmatrix}r_1\\r_2\end{pmatrix}
 =-\begin{pmatrix}X_{b+1}&Y_{b+1}\\X_{b+2}&Y_{b+2}\end{pmatrix}^{-1}
       \begin{pmatrix}\eta_{b+1}\\\eta_{b+2}\end{pmatrix}.
\end{equation}
Then $r_s=X_sr_1+Y_sr_2+\eta_s$ and
$Z_k=(-1)^{b-k}(2b-2k)!r_{b-k}$. This uses one two-by-two inverse
with a proved nonzero determinant. Equations \eqref{eg:backward} and
\eqref{eg:reconstruction} reconstruct all the even-zeta coordinates
from the first $q$ moments; substitution into \eqref{eg:entries} and
\eqref{eg:forward} recovers every later moment.

\subsection{An infinite family of ordinary rational polygamma sums}

Write the first row of $M_{p,q}^{-1}$ as
$(\lambda_b,\ldots,\lambda_{q-1})$ and define
\begin{equation}\label{eg:ordinary-weights}
 w_r=\sum_{j=\max\{b,r\}}^{q-1}
          \lambda_j\binom{2j+1}{2r}2^{-2(j-r)},\qquad
 P_b(x)=\sum_{r=0}^{q-1}w_rx^{2r}.
\end{equation}
For $p=q-4$, put
\begin{align*}
 u_0&=\frac1{(p-1)(q-1)},\\
 u_1&=-\frac1{24}\left(\frac q{p-1}+\frac p{q-1}\right),\\
 u_2&=\frac7{5760}\left(
 \frac{q(q+1)(q+2)}{p-1}+\frac{p(p+1)(p+2)}{q-1}\right)
          +\frac{pq}{576},\\
 S_b(x)&=u_0w_{q-1}x^4
       +(u_0w_{q-2}+u_1w_{q-1})x^2
       +u_0w_{q-3}+u_1w_{q-2}+u_2w_{q-1}.
\end{align*}

\begin{corollary}[Gap-four rational evaluations]\label{eg:ordinary-family}
For every even $q\ge6$, with $p=q-4$ and $x_n=n+1/2$,
\begin{equation}\label{eg:ordinary-family-equation}
 \boxed{\displaystyle
 \sum_{n=0}^{\infty}\left\{
 \frac{P_b(x_n)}{(p-1)!(q-1)!}
   \psi^{(p-1)}(n+1)\psi^{(q-1)}(n+1)-S_b(x_n)
                  \right\}=-\frac12.}
\end{equation}
The sum converges absolutely. Its polynomial $P_b$ has exact degree
$2q-2$, because
\begin{equation}\label{eg:top-coefficient}
 w_{q-1}=\frac{4b(4b-1)E_b}{d_{b+2}E_{b+1}}>0.
\end{equation}
\end{corollary}

\begin{proof}
The inverse row gives $\sum_{j=b}^{q-1}\lambda_jC_j=Z_0=-1/2$.
The inverse Bernoulli transform gives the weights
\eqref{eg:ordinary-weights}. The centered expansion of the tail product
is
\[
 A_p(n)A_q(n)=u_0x_n^{2-p-q}+u_1x_n^{-p-q}
                         +u_2x_n^{-p-q-2}+O(x_n^{-p-q-4}).
\]
Since $\deg P_b\le2q-2$ and $p=q-4$, its nonnegative powers are
exactly $S_b(x_n)$; the remaining summand is $O(x_n^{-2})$.
The continued centered sums of $x_n^4,x_n^2,1$ are all zero,
so the ordinary sum equals the moment combination. Equivalently,
finite centered power summation gives these polynomial subtraction
terms zero constant coefficient at the integer cutoff. The polygamma
normalization follows from even $p,q$.

For \eqref{eg:top-coefficient}, set every row output in
\eqref{eg:coefficient-equation} except the last equal to zero.
A homogeneous solution before the final boundary equation has initial
vector proportional to $(Y_{b+1},-X_{b+1})$, since its coefficient
at $b+1$ must vanish. Its coefficient at $b$ is $-E_b$ times that
scalar, and its continued coefficient at $b+2$ is $E_{b+1}$ times
that scalar. Hence the last row gives
\[
        \frac{r_b}{Y_b^*}=\frac{E_b}{d_{b+2}E_{b+1}}.
\]
The row and column scalings above show that
$\lambda_{q-1}=4bE_b/(d_{b+2}E_{b+1})$. Multiplication by
$\binom{2q-1}{2q-2}=4b-1$ gives the stated positive coefficient.
\end{proof}

\subsection{A concrete evaluation at the new pair \texorpdfstring{$(8,12)$}{(8,12)}}

A centered difference keeps the integer coefficients readable. Put
\begin{align*}
 K(y)=y^{13}\bigl(&469839+1269227y^2-11819495y^4\\
                 &+56443363y^6+83672540y^8+12834702y^{10}\bigr),\\
 P(x)&=K(x+\tfrac12)-K(x-\tfrac12),\\
 S(x)&=\frac{2951981460x^4+51273021735x^2+118231913818}{770}.
\end{align*}
Then the exact ordinary identity is
\begin{equation}\label{eg:8-12-sum}
 \boxed{\displaystyle
 \sum_{n=0}^{\infty}\left\{
 \frac{P(n+\tfrac12)}{7!\,11!}
       \psi^{(7)}(n+1)\psi^{(11)}(n+1)-S(n+\tfrac12)\right\}
 =-\frac{790768509}{16}.}
\end{equation}
For example, the coefficient row
\[
 (469839,1269227,-11819495,56443363,83672540,12834702)
\]
multiplied by $M_{8,12}$ is exactly
$(790768509/8,0,0,0,0,0)$. The centered difference converts this row
of blocks into $P$, and the first three centered tail-product
coefficients give precisely $S$. Thus \eqref{eg:8-12-sum} follows
already from the preceding corollary.

There is also a finite telescoping proof, independent of the
continued-moment interpretation. Let $c_0,\ldots,c_5$ be the six
coefficients in the parentheses defining $K$, in increasing order,
and let
\[
 F_m(N)=\sum_{k=1}^N k^m
       =\frac{B_{m+1}(N+1)-B_{m+1}}{m+1}.
\]
Define the explicitly specified polynomials
\begin{align*}
 U_8(N)&=\sum_{j=0}^5c_jF_{2j+1}(N),&
 U_{12}(N)&=\sum_{j=0}^5c_jF_{2j+5}(N),\\
 V(N)&=-\frac{N}{6160}\bigl(
 2417202210N^4-5764920315N^3+42395905270N^2\notag\\
 &\hspace{36mm}-129829345800N+260578219219\bigr).
\end{align*}
The first two definitions imply
\[
 \Delta U_8(N)=K(N)N^{-12},\qquad
 \Delta U_{12}(N)=K(N)N^{-8},
\]
where $\Delta f(N)=f(N)-f(N-1)$. Direct polynomial expansion also
gives
\begin{equation}\label{eg:V-difference}
 \Delta V(N)=\frac{K(N)}{N^{20}}+
      \frac{U_8(N-1)}{N^8}+\frac{U_{12}(N-1)}{N^{12}}
                         -S(N-\tfrac12).
\end{equation}
In particular the negative powers cancel on the right. These are
finite rational-function identities, reproducible from the six
integer coefficients and the displayed Bernoulli polynomials.
Now put
\[
 \mathcal F(N)=K(N)A_8(N)A_{12}(N)
                 +U_8(N)A_8(N)+U_{12}(N)A_{12}(N)+V(N).
\]
Using $A_m(N-1)=A_m(N)+N^{-m}$ and the three difference identities,
all mixed terms cancel to give
\[
 \mathcal F(N)-\mathcal F(N-1)
 =P(N-\tfrac12)A_8(N-1)A_{12}(N-1)-S(N-\tfrac12).
\]
Since $\mathcal F(0)=0$, the partial sum of
\eqref{eg:8-12-sum} through $n=N-1$ is exactly $\mathcal F(N)$.
For the limiting step, use
\begin{align*}
 A_8(N)&=\frac1{7N^7}-\frac1{2N^8}+\frac2{3N^9}
                                      -\frac1{N^{11}}+O(N^{-13}),\\
 A_{12}(N)&=\frac1{11N^{11}}-\frac1{2N^{12}}+\frac1{N^{13}}
                                 -\frac{91}{30N^{15}}+O(N^{-17}).
\end{align*}
The three tail-containing terms of $\mathcal F$ have constant
coefficients, respectively,
\[
       -\frac{116518888243}{2310},\qquad
       -\frac{193499245}{336},\qquad
        \frac{350665911}{220}.
\]
Their sum is $-790768509/16$, $V$ has zero constant coefficient,
and the coefficients of $N,N^2,\ldots,N^5$ cancel exactly.
The remainders above contribute $O(N^{-1})$. This proves the
ordinary evaluation by a finite identity and an asymptotic limit.

\subsection{Scope and reproducible checks}

The exact script checks \eqref{eg:determinant} against the original
Bernoulli matrix for $3\le b\le18$, checks the independent rotation
matrix, and verifies \eqref{eg:top-coefficient} where a full inverse
is computed. It also checks 1,946 minors of order two for gaps
$c=2,\ldots,6$, and verifies the displayed ordinary evaluation by
an exact finite telescoper. These checks detect normalization errors;
the all-order proof is Theorem~\ref{eg:two-minors}, not the finite
arithmetic range.

The theorem about minors of order two cannot be promoted to a
fixed-sign assertion for larger boundary determinants. For example,
for $c=3$ the exact signs are
\[
 \det(R_{6+i,j})_{1\le i,j\le3}<0,\qquad
 \det(R_{7+i,j})_{1\le i,j\le3}>0.
\]
Both are nonzero, and both arrays satisfy all the strict negative
minor statements in Theorem~\ref{eg:two-minors}. Thus order-two
comparison alone does not prove the general conjecture. Combining
all currently established families, the still-unsettled range is
\[
             8\le p<q,\qquad q-p\ge6,
             \qquad p,q\ \hbox{even}.
\]
A useful next problem is to identify an exact, sign-flexible
nonvanishing invariant for the higher boundary determinants, or to
combine the fixed-gap recurrence with the complementary determinant
of size $p/2$ from the shifted-Hankel approach. The rank problem and
the arithmetic independence of the actual zeta coordinates remain
distinct questions.

% Contribution: Hurwitz-shifted tail moments and exact discrete primitives.
% Required bibliography keys: DLMF, DLMFBernoulli.
% All statements below were derived independently in the current continuation.
% Historical priority over the full literature is not asserted.
\section{Hurwitz-shifted tails: Bernoulli primitives and exact corrections}
\label{sh:section}

The centered-tail results in the inspected reports use the lattice
$n+1/2$ and the tails $\zeta(p,n+1)$.  Their ordinary subtraction rule
uses the special vanishing $\zeta(-2j,1/2)=0$.  A shift changes both
the lower endpoint of the finite summation and the continued sums of
the subtracted polynomials.  The following formulas retain both changes.
They answer the exact-reduction part of the proposed Hurwitz-shifted
extension; they do not assert a new determinant nonvanishing theorem.

The ingredients are the classical Hurwitz recurrence, the Bernoulli
difference law, and Euler--Maclaurin continuation
\cite{DLMF,DLMFBernoulli}.  The particular discrete
primitive, shifted block formulas, and normalized identities are proved
below without a numerical relation search.

\subsection{An explicit discrete primitive at every sufficiently large degree}

For an integer $s$, use the pole-free convention
\[
 \mathcal Z_s(v)=
 \begin{cases}
  \zeta(s,v),&s\ne1,\\
  -\psi(v),&s=1.
 \end{cases}
\]
Thus $\mathcal Z_s(v)-\mathcal Z_s(v+1)=v^{-s}$ at every
integer $s$, including $s=1$.  For integers $s\ge2$ and $m\ge0$,
define the finite expression
\begin{equation}\label{sh:L}
 \mathcal L_{s,m}(v)=\frac1{m+1}
    \sum_{\substack{0\le h\le m+1\\h\text{ even}}}
       \binom{m+1}{h}B_h\,
       \mathcal Z_{s-m-1+h}(v)
       -\frac{B_{m+1}(v)}{m+1}\zeta(s,v).
\end{equation}
All nonpositive integer Hurwitz values in this formula are Bernoulli
polynomials.  In particular the formula is finite even when $m$ is large.

\begin{theorem}[Bernoulli discrete primitive]\label{sh:primitive}
Let $p,q\ge2$ be integers, let $d\ge\max(p,q)$ be an integer,
and let $v>0$.  Put
\begin{align}
 \mathcal C_d^{p,q}(v)
   &=\mathcal L_{p,d-q}(v)+\mathcal L_{q,d-p}(v),\label{sh:Chigh}\\
 \mathcal K_d^{p,q}(v)
   &=\mathcal C_d^{p,q}(v)
       -(v-1)^d\zeta(p,v)\zeta(q,v).\label{sh:K}
\end{align}
Then
\begin{equation}\label{sh:antidifference}
 \boxed{\mathcal K_d^{p,q}(v)-\mathcal K_d^{p,q}(v+1)
  =\{v^d-(v-1)^d\}\zeta(p,v)\zeta(q,v).}
\end{equation}
Consequently, for every integer $N\ge1$,
\begin{equation}\label{sh:finitetelescoping}
 \sum_{n=0}^{N-1}
 \{(n+v)^d-(n+v-1)^d\}\zeta(p,n+v)\zeta(q,n+v)
 =\mathcal K_d^{p,q}(v)-\mathcal K_d^{p,q}(v+N).
\end{equation}
If also $d<p+q-2$, the infinite series is ordinary and absolutely
convergent, and its value is $\mathcal K_d^{p,q}(v)$.
\end{theorem}

\begin{proof}
Write $R_m(v)=B_{m+1}(v)/(m+1)$.  The Bernoulli difference law is
$R_m(v+1)-R_m(v)=v^m$.  Because $B_1=-1/2$ and all higher odd
Bernoulli numbers vanish,
\[
 \frac1{m+1}\sum_{h\text{ even}}
       \binom{m+1}{h}B_hv^{m+1-h}
 =R_m(v)+\frac12v^m.
\]
Use these two identities together with the Hurwitz recurrence in
\eqref{sh:L}.  The result is exactly
\begin{equation}\label{sh:Ldifference}
 \mathcal L_{s,m}(v)-\mathcal L_{s,m}(v+1)
   =v^m\zeta(s,v)-\frac12v^{m-s}.
\end{equation}
Thus $\mathcal C_d(v)-\mathcal C_d(v+1)$ equals
\[
 v^{d-q}\zeta(p,v)+v^{d-p}\zeta(q,v)-v^{d-p-q}
 =v^d\{\zeta(p,v)\zeta(q,v)
              -\zeta(p,v+1)\zeta(q,v+1)\}.
\]
Subtracting the difference of $(v-1)^d\zeta(p,v)\zeta(q,v)$
gives \eqref{sh:antidifference}; summing gives
\eqref{sh:finitetelescoping}.  In the stated convergent range,
the summand is $O(n^{d+1-p-q})$.  Direct use of the Hurwitz
asymptotic expansion in \eqref{sh:K} gives
$\mathcal K_d^{p,q}(v)=O(v^{d+2-p-q})$, hence a zero limit at
infinity.  Alternatively this last bound follows term by term in
\eqref{sh:L}, since all its Hurwitz orders are then at least two.
\end{proof}

The convention at $s=1$ in \eqref{sh:L} is essential.  Inserting
$\zeta(1,v)$ into the displayed finite expression would make a valid
discrete primitive appear divergent.  Its correct coordinate is
$\mathcal Z_1(v)=\gamma_0(v)=-\psi(v)$.

\subsection{Every even centered moment, with its lower endpoint}

Assume now that $2\le p\le q$ are even and $v>1/2$.  Put
\[
 \alpha=v-\frac12,\qquad a=v-1,\qquad x_n=n+\alpha,
 \qquad A_p(n;v)=\zeta(p,n+v),
\]
and let
\begin{equation}\label{sh:D}
 D_{p,q}(s;v)=\sum_{n\ge0}
           A_p(n;v)A_q(n;v)x_n^{-s},
 \qquad \Re s>3-p-q.
\end{equation}
Its meromorphic continuation is understood outside this half-plane.
For $r\ge0$ set $D_r(v)=D_{p,q}(-2r;v)$ and
\begin{equation}\label{sh:Dhat}
 \widehat D_r(v)=D_r(v)
       -\zeta(-2r,\alpha)\zeta(p,v)\zeta(q,v).
\end{equation}
For $u\ge2$, $w\ge1$, use the convergent strict double Hurwitz value
\[
 Z_v(u,w)=\sum_{\ell>k\ge0}
                  (\ell+v)^{-u}(k+v)^{-w}.
\]

\begin{theorem}[Complete shifted Bernoulli transform]\label{sh:transform}
Every $D_r(v)$ is regular.  Define $\mathcal C_d(v)$ for positive
odd $d$ by the following three formulas:
\begin{align}
 \mathcal C_d(v)
 &=Z_v(p,q-d)+Z_v(q,p-d)+\zeta(p+q-d,v),
 &&d<p,\label{sh:Clow}\\
 \mathcal C_d(v)
 &=Z_v(p,q-d)+\frac12\zeta(p+q-d,v)
                         +\mathcal L_{q,d-p}(v),
 &&p<d<q,\label{sh:Cmiddle}\\
 \mathcal C_d(v)
 &=\mathcal L_{p,d-q}(v)+\mathcal L_{q,d-p}(v),
 &&d>q.\label{sh:Cupper}
\end{align}
The middle range is empty when $p=q$.  Since $p,q$ are even and
$d$ is odd, these three ranges cover all required indices.  Then
\begin{align}
 \widehat D_r(v)
 &=\sum_{j=0}^r
    \frac{\binom{2r}{2j}}{2j+1}
       B_{2r-2j}(1/2)\,\mathcal C_{2j+1}(v),
       \label{sh:forward}\\
 \mathcal C_{2j+1}(v)
 &=\sum_{r=0}^j\binom{2j+1}{2r}
                         2^{-2(j-r)}\widehat D_r(v).
       \label{sh:inverse}
\end{align}
At $v=1$ the correction in \eqref{sh:Dhat} vanishes.  For a
general shift it is a specified Bernoulli polynomial multiplied by
the initial product of tails.
\end{theorem}

\begin{proof}
The centered expansion is
\begin{equation}\label{sh:centeredasymptotic}
 \zeta(p,x+1/2)\sim x^{1-p}
 \left\{\frac1{p-1}+\sum_{h\ge1}
   \frac{B_{2h}(1/2)(p)_{2h-1}}{(2h)!}x^{-2h}\right\}.
\end{equation}
Consequently the only possible poles of \eqref{sh:D} are at
$s=3-p-q-2h$, all of them odd integers.  Finite subtraction of
\eqref{sh:centeredasymptotic} proves regularity at every $-2r$.

Let $Q_r(y)=B_{2r+1}(y+1/2)/(2r+1)$.  It is odd, and
\[
 \sum_{n=0}^{N-1}x_n^{2r}=Q_r(N+a)-Q_r(a),\qquad
 Q_r(a)=-\zeta(-2r,\alpha).
\]
Finite summation by parts with
$F(t)=\zeta(p,t)\zeta(q,t)$ therefore expresses the partial
moment as
\begin{equation}\label{sh:finiteEtransform}
 \sum_{j=0}^r
   \frac{\binom{2r}{2j}}{2j+1}B_{2r-2j}(1/2)
        E_{2j+1}(N;v)-Q_r(a)F(v),
\end{equation}
where, with $y=N+v-1$,
\begin{align*}
 E_d(N;v)={}&y^dF(y+1)
  +\sum_{k=0}^{N-1}(k+v)^{d-p}\zeta(q,k+v+1)\\
 &+\sum_{k=0}^{N-1}(k+v)^{d-q}\zeta(p,k+v+1)
  +\sum_{k=0}^{N-1}(k+v)^{d-p-q}.
\end{align*}
All equalities so far are finite identities.

Take constant terms as $y\to\infty$.  The continued value of the
partial moment equals its constant term in this coordinate: every
nonnegative asymptotic monomial is an even power of $x$, and
$\zeta(-2h,N+\alpha)=\zeta(-2h,y+1/2)$ is an odd polynomial in
$y$, hence contributes zero constant term at the upper endpoint.

For $d<p$ every sum in $E_d$ converges and the boundary tends to
zero, giving \eqref{sh:Clow}.  For $p<d<q$, apply the finite
Faulhaber identity
\begin{align*}
 \sum_{k=0}^{N-1}(k+v)^m\zeta(s,k+v+1)
 ={}&\frac{B_{m+1}(y+1)-B_{m+1}(v)}{m+1}\zeta(s,y+1)\\
 &+\frac1{m+1}\sum_{k=0}^{N-1}
      \frac{B_{m+1}(k+v)-B_{m+1}(v)}{(k+v)^s}
\end{align*}
to its one nonnegative power.  Here $m=d-p$ is odd and
$m<q-1$, so no endpoint term remains.  Its $B_1=-1/2$ term
combines with the final sum in $E_d$ to give
$\tfrac12\zeta(p+q-d,v)$; the other terms are precisely
\eqref{sh:Cmiddle}.

For $d>q$, Theorem~\ref{sh:primitive} gives the exact identity
$E_d(N;v)=\mathcal C_d(v)-\mathcal K_d(y+1)$.
For completeness the last endpoint has zero constant term.
Write $c_{s,j}=[y^{-j}]\zeta(s,y+1)$.  For even $s$, every
exponent in this expansion is odd except the term
$-\tfrac12y^{-s}$.  Since $m$ is odd, direct constant-term
extraction in \eqref{sh:L} gives
\[
 \operatorname{CT}_y\mathcal L_{s,m}(y+1)
 =-\frac12c_{s,m}.
\]
Indeed the constant from the first finite sum cancels the
even-power part of $B_{m+1}(y+1)\zeta(s,y+1)/(m+1)$; its
remaining contribution is $c_{s,m}/2$, from
$B_1(1)=1/2$.  On the other hand
\[
 \operatorname{CT}_y\{y^dF(y+1)\}
 =-\frac12\{c_{p,d-q}+c_{q,d-p}\}.
\]
Subtracting proves $\operatorname{CT}_y\mathcal K_d(y+1)=0$.
This also covers $d=p+q-1$, where individual harmonic logarithms
in finite Faulhaber sums cancel only after all terms are combined.
Thus \eqref{sh:Cupper} holds at every required high index.

Taking constant terms in \eqref{sh:finiteEtransform} proves
\eqref{sh:forward}.  Finally the Bernoulli polynomial identity
\[
 y^{2j+1}=\sum_{r=0}^j
       \binom{2j+1}{2r}2^{-2(j-r)}Q_r(y)
\]
inverts the finite transform and gives \eqref{sh:inverse}.
\end{proof}

For $d>q$, every order in \eqref{sh:Cupper} is even and no larger
than $q$.  Orders below zero are Bernoulli polynomials.  It follows
that, over $\mathbb Q(v)$, all even moments lie in the span of
\[
 \mathcal C_1(v),\mathcal C_3(v),\ldots,\mathcal C_{q-1}(v),
 \quad 1,\zeta(2,v),\zeta(4,v),\ldots,\zeta(q,v),
 \quad\zeta(p,v)\zeta(q,v).
\]
This is a list of $q+2$ specified coordinates, not an assertion of
independence.  At $v=1$ the initial-product correction vanishes and
the highest single-zeta column cancels, recovering the smaller
unshifted coordinate system.  At rational $v$, all displayed
coefficients and all negative even Hurwitz values are rational.

\subsection{Ordinary sums and every Stieltjes correction}

The following statement also covers odd tail orders and odd polynomial
weights.  Its derivatives are in the Dirichlet weight $x_n^{-s}$,
not in the spectral orders $p,q$ of the two tails.

\begin{theorem}[Exact subtraction and all weight jets]\label{sh:jets}
Let $p,q\ge2$ be integers, $v>1/2$, $P$ a polynomial, and
\[
 f_n=P(x_n)\zeta(p,n+v)\zeta(q,n+v).
\]
Let $R(X)=\sum_{k\ge0}r_kX^k$ be the nonnegative-power part of
its asymptotic expansion in $x_n$, and let $\rho$ be the coefficient
of $x_n^{-1}$.  These are finite coefficient calculations from
\eqref{sh:centeredasymptotic}.  Set
$u_n=f_n-R(x_n)-\rho/x_n$.  Then $u_n=O(x_n^{-2})$ and
\begin{equation}\label{sh:jetmaster}
 \sum_{n\ge0}f_nx_n^{-s}
 =\sum_{n\ge0}u_nx_n^{-s}
       +\sum_{k\ge0}r_k\zeta(s-k,\alpha)
       +\rho\zeta(1+s,\alpha)
\end{equation}
is a meromorphic identity near $s=0$, with residue $\rho$.  For
every $j\ge0$,
\begin{align}\label{sh:alljets}
 [s^j]\left(\sum_{n\ge0}f_nx_n^{-s}-\frac\rho s\right)
 ={}&\frac{(-1)^j}{j!}\sum_{n\ge0}u_n\log^j x_n
     +\frac1{j!}\sum_{k\ge0}r_k\zeta^{(j)}(-k,\alpha)\notag\\
 &+\frac{(-1)^j\rho}{j!}\gamma_j(\alpha).
\end{align}
All sums on the right are ordinary and absolutely convergent.
In particular, when $p,q$ and $P$ are even, $\rho=0$ and
\begin{equation}\label{sh:ordinarycorrection}
 \boxed{\sum_{n\ge0}\{f_n-R(x_n)\}
   =\sum_r[X^{2r}]P(X)\,D_r(v)
                       -\sum_{k\ge0}r_k\zeta(-k,\alpha).}
\end{equation}
\end{theorem}

\begin{proof}
Subtract the finitely many indicated terms in an initial right
half-plane of absolute convergence.  The remainder series is normally
convergent for $\Re s>-1$, as are all its fixed $s$ derivatives
on compact subsets.  This proves \eqref{sh:jetmaster} by unique
meromorphic continuation.  Expand the ordinary Hurwitz terms and
\[
 \zeta(1+s,\alpha)=\frac1s+
       \sum_{j\ge0}\frac{(-1)^j}{j!}\gamma_j(\alpha)s^j.
\]
Coefficient comparison gives \eqref{sh:alljets}.  The even case has
only even powers of $x_n$, so no $x_n^{-1}$ term occurs.  Its
constant coefficient is \eqref{sh:ordinarycorrection}.
\end{proof}

The polynomial correction in \eqref{sh:ordinarycorrection} vanishes
on the original lattice $\alpha=1/2$.  It must generally be retained
after a shift.  For example at $\alpha=1$, a subtracted constant
$c$ contributes $-c\zeta(0,1)=c/2$, even though all higher
negative even zeta values vanish there.

\subsection{Explicit trigamma identities and derivative families}

Let $T(v)=\psi'(v)=\zeta(2,v)$.  Formula \eqref{sh:Chigh} with
$p=q=2$, $d=3$ gives
\[
 \mathcal C_3(v)=\frac12-v+v(1-v)T(v),\qquad
 \mathcal K_3(v)=\frac12-v+v(1-v)T(v)-(v-1)^3T(v)^2.
\]
The summand has constant asymptotic part $3$.  Applying
\eqref{sh:ordinarycorrection} yields the ordinary identity
\begin{equation}\label{sh:trigammasquare}
 \boxed{\begin{aligned}
 &\sum_{n\ge0}\left[
    \left\{3(n+v-\tfrac12)^2+\tfrac14\right\}
                     \psi'(n+v)^2-3\right]\\
 &\qquad=2v-\frac52+v(1-v)\psi'(v)
                         -(v-1)^3\psi'(v)^2.
 \end{aligned}}
\end{equation}
Initially $v>1/2$; normal convergence extends this identity to every
real $v>0$, and indeed to $\Re v>0$.  Two exact specializations are
\begin{align}
 \sum_{n\ge0}\left[
       (3n^2+3n+1)(\zeta(2)-H_n^{(2)})^2-3\right]
       &=-\frac12,\label{sh:unshiftedsquare}\\
 \sum_{n\ge0}\left[
       \{3(n+1)^2+\tfrac14\}\zeta(2,n+\tfrac32)^2-3\right]
       &=\frac32+\frac{\pi^2}{8}-\frac{\pi^4}{32}.
       \label{sh:halflatticesquare}
\end{align}
The second follows from $\zeta(2,3/2)=\pi^2/2-4$.

All parameter derivatives are again ordinary convergent identities.
For example, with $x_n=n+v-1/2$,
\begin{align}\label{sh:squarederivative}
 &\sum_{n\ge0}\left[
  6x_n\psi'(n+v)^2
  +2(3x_n^2+\tfrac14)\psi'(n+v)\psi''(n+v)\right]\notag\\
 &\quad=2+(1-2v)\psi'(v)+v(1-v)\psi''(v)
      -3(v-1)^2\psi'(v)^2
      -2(v-1)^3\psi'(v)\psi''(v).
\end{align}
At $v=1$ this evaluates the displayed mixed polygamma sum to
$2-\zeta(2)$.  More generally, the $m$th derivative of the summand
in \eqref{sh:trigammasquare}, for $m\ge1$, is the explicit finite
Leibniz sum
\[
 \sum_{j=0}^{\min(m,2)}\binom mj P^{(j)}(x_n)
    \sum_{k=0}^{m-j}\binom{m-j}{k}
       \psi^{(k+1)}(n+v)\psi^{(m-j-k+1)}(n+v),
 \qquad P(X)=3X^2+\tfrac14.
\]
Its ordinary sum is the corresponding derivative of the right side
of \eqref{sh:trigammasquare}.  Normal convergence follows either
by differentiated Euler--Maclaurin remainders or by Cauchy estimates
on compact subsets of $\Re v>0$.

\begin{corollary}[Every derivative at the integral lattice]\label{sh:allderivatives}
Put $P_0(z)=3z^2-3z+1$ and, for $k\ge1$, define the ordinary sum
\[
 S_k=\frac1{k!}\sum_{n\ge0}
     \left.\frac{d^k}{dz^k}\{P_0(z)\psi'(z)^2\}
                                  \right|_{z=n+1}.
\]
Then $S_1=2-\zeta(2)$, and for every $k\ge2$,
\begin{equation}\label{sh:derivativeclosed}
 \boxed{S_k=(-1)^k\left\{
      k\zeta(k+1)-(k-1)\zeta(k)
      +\sum_{a=2}^{k-1}(a-1)(k-a)
                     \zeta(a)\zeta(k+1-a)\right\}.}
\end{equation}
The last sum is empty for $k=2$.  This gives finite evaluations of
an infinite family of mixed polygamma sums in ordinary zeta values
and products of two such values.
\end{corollary}

\begin{proof}
Write $h=v-1$ and $T=\psi'(1+h)$.  The right side of
\eqref{sh:trigammasquare} becomes
\[
 -\frac12+2h-(h+h^2)T-h^3T^2,
 \qquad
 T=\sum_{j\ge0}(-1)^j(j+1)\zeta(j+2)h^j,
 \quad |h|<1.
\]
The coefficient of $h$ is $2-\zeta(2)$.  For $k\ge2$, the
coefficients from the two linear terms in $T$ are
$(-1)^k\{k\zeta(k+1)-(k-1)\zeta(k)\}$; the coefficient
from $-h^3T^2$ is the finite convolution in
\eqref{sh:derivativeclosed}.  Normal convergence, proved above,
identifies these coefficients with the stated sums.
\end{proof}

There is also a simple resonant example.  Taking $p=q=d=2$ in
Theorem~\ref{sh:primitive} gives
\[
 \mathcal K_2(v)=-2\psi(v)+(1-2v)T(v)-(v-1)^2T(v)^2.
\]
Here the summand has leading term $2/x_n$, and
$\mathcal K_2(v+N)=-2\log N-3+o(1)$.  Comparing this expansion
with the finite harmonic sum proves
\begin{equation}\label{sh:resonantordinary}
 \boxed{\begin{aligned}
 &\sum_{n\ge0}\left[
     (2n+2v-1)\psi'(n+v)^2-\frac2{n+v-1/2}\right]\\
 &\quad=3-2\psi(v)+2\psi(v-\tfrac12)
                 +(1-2v)\psi'(v)-(v-1)^2\psi'(v)^2,
 \qquad v>\tfrac12.
 \end{aligned}}
\end{equation}
Its Dirichlet finite part is instead
\begin{equation}\label{sh:resonantFP}
 \operatorname{FP}_{s=0}
   \sum_{n\ge0}(2n+2v-1)\psi'(n+v)^2x_n^{-s}
 =3-2\psi(v)+(1-2v)\psi'(v)-(v-1)^2\psi'(v)^2.
\end{equation}
The difference is exactly $2\gamma_0(v-1/2)$ as prescribed by
Theorem~\ref{sh:jets}.  In particular the ordinary sum in
\eqref{sh:resonantordinary} at $v=1$ is
$3-4\log2-\zeta(2)$, whereas its continued finite part is
$3+2\gamma-\zeta(2)$.

For a convergent higher-order example requiring no subtraction,
$p=q=4$, $d=5$ gives
\begin{equation}\label{sh:fourthsquare}
 \sum_{n\ge0}\{(n+v)^5-(n+v-1)^5\}\zeta(4,n+v)^2
 =\zeta(2,v)+v(1-v)\zeta(4,v)
                   -(v-1)^5\zeta(4,v)^2,
 \qquad v>0.
\end{equation}
These formulas are exact finite reductions in single Hurwitz values
and endpoint products.  They do not require treating a nonzero formal
rank or a small numerical residual as an arithmetic independence claim.

\section{An alternating two-parameter digamma identity}
\label{alt:sec:master}

The centered moments in the preceding sections use even polynomial weights.
An alternating outer sign has a complementary effect: odd centered weights
lead to odd-weight alternating double zeta values, which admit a classical
parity reduction. Summing these reductions over both even tail orders gives
a compact identity in digamma and trigonometric functions. The method of
turning products of tails into ordered sums is classical; see Hoffman
\cite{Hoffman}. The parity formula is due to the classical Euler-sum
theory and is proved in a generating-function form by Teo
\cite{Teo}; see also the descending-index formulation in
\cite[Eq.~(3.7)]{XuWang}. The contribution here is the resulting
two-parameter identity, its explicitly normalized odd-weight family, and
the rational-parameter square evaluations. No literature-priority claim
is made for individual special cases.

For $\sigma\in\{1,-1\}$ set
\begin{equation}\label{alt:eq:resolvent-tail}
 E_n^\sigma(x)=\sum_{k>n}\frac{\sigma^k}{k^2-x^2},
 \qquad x\notin\mathbb Z\setminus\{0\},\quad n\ge0.
\end{equation}
Each $E_n^\sigma$ is even in $x$. Its poles at integers of modulus at
most $n$ are absent. These functions are ordinary polygamma expressions:
\begin{align}
 E_n^+(x)&=\frac{\psi(n+1+x)-\psi(n+1-x)}{2x},
 \label{alt:eq:digamma-tail}\\
 E_n^-(x)&=\frac{(-1)^{n+1}}{2x}
       \{\mathcal B(n+1-x)-\mathcal B(n+1+x)\},
 \qquad
 \mathcal B(v)=\frac{\psi((v+1)/2)-\psi(v/2)}2.
 \notag
\end{align}
The apparent singularities at $x=0$ are removable. Formula
\eqref{alt:eq:digamma-tail} follows by partial fractions and the
convergent digamma difference series; the second formula follows from
$\mathcal B(v)=\sum_{j\ge0}(-1)^j/(j+v)$ for $\Re v>0$,
and then by continuation.

Define the even meromorphic functions
\begin{align}
 O_+(t)&=-\frac{\psi(1+t)+\psi(1-t)+2\gamma}{2},
 \label{alt:eq:O}\\
 O_-(t)&=-\log2+O_+(t/2)-O_+(t),\notag
\end{align}
and put $R_+(t)=\cot(\pi t)$ and $R_-(t)=\csc(\pi t)$.
For $xy\ne0$, abbreviate
\begin{equation}\label{alt:eq:UV}
 U_\sigma(x,y)=\frac{O_\sigma(x+y)-O_\sigma(x-y)}{2xy},
 \qquad
 V_\sigma(x,y)=\frac{O_\sigma(x+y)+O_\sigma(x-y)}2.
\end{equation}

\begin{theorem}[The two-parameter alternating identity]
\label{alt:thm:master}
Let $\alpha,\beta\in\{1,-1\}$ and $c=-\alpha\beta$. Initially assume
$|x|+|y|<1$ and $xy\ne0$. Then
\begin{align}
 &\mathcal K_{\alpha,\beta}(x,y)
 :=\sum_{n=0}^\infty(-1)^n(2n+1)E_n^\alpha(x)E_n^\beta(y)
 \notag\\
 &=\frac\pi2\bigg\{
 \frac{R_\alpha(x)}xO_{-\beta}(y)
 +\frac{R_\beta(y)}yO_{-\alpha}(x)
 -yR_c(y)U_\alpha(x,y)-xR_c(x)U_\beta(x,y)
 \notag\\[-2pt]
 &\hspace{37mm}
 -\frac{R_c(x)}xV_{-\beta}(x,y)
 -\frac{R_c(y)}yV_{-\alpha}(x,y)\bigg\}.
 \label{alt:eq:master}
\end{align}
The series converges absolutely and locally normally if at least one of
$\alpha,\beta$ is $-1$. For $\alpha=\beta=1$ it converges ordinarily
and locally uniformly; after subtracting each summand's value at
$(0,0)$ the remaining series is absolutely locally normally
convergent. Both sides extend meromorphically to $\mathbb C^2$.
Possible genuine poles occur only when $x$ or $y$ is a nonzero integer.
In particular the apparent singularities at $xy=0$, and those on
$x\pm y\in\mathbb Z$ away from the coordinate poles, are removable.
The identity is even separately in $x$ and $y$ and is unchanged by
simultaneously exchanging $(\alpha,x)$ with $(\beta,y)$.
\end{theorem}

\subsection{An ordinary finite-cutoff identity}

For $s\ge2$, write
\[
 T_s^\sigma(n)=\sum_{k>n}\frac{\sigma^k}{k^s},\qquad
 Z(r,s;\sigma,\tau)
 =\sum_{k>m\ge1}\frac{\sigma^k\tau^m}{k^r m^s}.
\]
The convention for $Z$ is the decreasing-index convention. Put
$\mathsf L_s(\sigma)=\operatorname{Li}_s(\sigma)$ for $s\ge2$ and
$\mathsf L_1(-1)=-\log2$.

\begin{lemma}[Odd alternating weights]
\label{alt:lem:finitecutoff}
Let $p,q\ge2$ be integers and let $d$ be odd with
$1\le d<\min(p,q)$. Then
\begin{align}
 &\sum_{n\ge0}(-1)^n\{n^d+(n+1)^d\}
                    T_p^\alpha(n)T_q^\beta(n)\notag\\
 &\quad=-Z(p,q-d;\alpha,-\beta)
        -Z(q,p-d;\beta,-\alpha)
        -\mathsf L_{p+q-d}(-\alpha\beta).
 \label{alt:eq:weighted-tail}
\end{align}
The left side is absolutely convergent except possibly when
$(p,q,d,\alpha,\beta)=(2,2,1,1,1)$; in that case it converges as
an ordinary alternating series. Every double series on the right is
absolutely convergent.
\end{lemma}

\begin{proof}
For $m\ge1$, cancellation of consecutive powers gives
\[
 \sum_{n=0}^{m-1}(-1)^n\{n^d+(n+1)^d\}
       =(-1)^{m-1}m^d.
\]
First sum over the finite range $0\le n\le N$. For each fixed $n$
the two tail series are absolutely convergent, so their product may
be expanded and the finite $n$-sum interchanged. The inner weight
becomes $(-1)^{\min(k,m,N+1)-1}\min(k,m,N+1)^d$.
The limiting double series is dominated by
$\min(k,m)^d/(k^p m^q)$, whose sum is finite because $q-d,p-d\ge1$
and $p,q\ge2$. For example the region $k>m$ is
\[
 \sum_{m\ge1}m^{d-q}\sum_{k>m}k^{-p}
       =O\left(\sum_{m\ge1}m^{d+1-p-q}\right)<\infty.
\]
Dominated convergence is therefore applicable to the double series.
Separating $k>m$, $m>k$, and $k=m$ proves
\eqref{alt:eq:weighted-tail}. Equivalently, the additional term in the
finite-cutoff formula when both indices exceed $N+1$ has modulus at
most $(N+1)^d T_p^+(N)T_q^+(N)=O(N^{d+2-p-q})$, which tends to zero.

For absolute convergence of the original series use
$T_s^+(n)=O(n^{1-s})$ and $T_s^-(n)=O(n^{-s})$.
Under the stated range the only unresolved exponent is $-1$ in the
listed case. There
$(2n+1)T_2^+(n)^2=2/n+O(n^{-2})$ for $n\ge1$, so an ordinary
alternating harmonic series plus an absolutely convergent series
proves convergence. This also fixes the boundary value without
assigning a value to a divergent series.
\end{proof}

\subsection{The classical parity input, with its endpoint convention}

The following formula is the alternating double-Euler parity theorem
\cite{Teo,XuWang}. We record it in the convention used here.
For $r\ge2$, $s\ge1$, and $w=r+s$ odd, set
\[
 \mu_j=(-1)^r\left\{
 \binom{j-1}{r-1}\mathsf L_j(\sigma)
 +\binom{j-1}{s-1}\mathsf L_j(\tau)\right\}.
\]
Binomial coefficients outside $0\le k\le n$ are zero. Then
\begin{align}
 Z(r,s;\sigma,\tau)
 ={}&\frac12\left\{-\mathsf L_w(\sigma\tau)
 +(1+(-1)^r)\mathsf L_r(\sigma)\mathsf L_s(\tau)+\mu_w\right\}
 \notag\\
 &-\sum_{k=1}^{(w-1)/2}\mathsf L_{2k}(\sigma\tau)\mu_{w-2k}.
 \label{alt:eq:parity}
\end{align}
The exceptional notation $\mathsf L_1(1)$ is interpreted as one formal
symbol $T$ while this finite expression is assembled. It cancels before
the expression is evaluated. Indeed it can occur only when $s=1$ and
$\tau=1$; then $r$ is even, and the first product contributes
$\mathsf L_r(\sigma)T$, while the last summand contributes its negative.
There is no remaining regularization parameter. Thus one may set
$\mathsf L_1(1)=0$ in a coefficient implementation after this cancellation
has been established. The value zero is a bookkeeping convention, not
an asserted value of the divergent harmonic series.

Combining \eqref{alt:eq:weighted-tail} and \eqref{alt:eq:parity}
already proves an explicit all-order family whenever $p+q$ is even.
Only single zeta values and $\log2$ remain, since
$\mathsf L_s(-1)=(2^{1-s}-1)\zeta(s)$ for $s\ge2$.
No period-independence assumption is involved.
The canonical weights exhaust the corresponding odd polynomial space:
for odd $d$,
$P_d(X)=(X-1/2)^d+(X+1/2)^d$ is odd with leading term $2X^d$.
Consequently the $P_d$ with $d<\min(p,q)$ form a triangular basis
for all odd centered polynomial weights in that degree range.
For $p+q$ even, every such weighted moment therefore has the same
finite reduction to single zeta values and $\log2$.

\subsection{Proof of the two-parameter identity}

Put
\[
 A_\sigma(t)=\sum_{a\ge1}\mathsf L_{2a}(\sigma)t^{2a-2}
       =\frac1{2t^2}-\frac{\pi R_\sigma(t)}{2t}.
\]
The equality follows from the cotangent partial-fraction expansion and
its alternating counterpart. Moreover, for $|t|<1$,
\[
 O_\sigma(t)=\sum_{j\ge0}\mathsf L_{2j+1}(\sigma)t^{2j},
\]
with the already justified convention $\mathsf L_1(1)=0$.
The familiar digamma Taylor expansion gives \eqref{alt:eq:O}.

We first verify the analytic legitimacy of the coefficient calculation.
Uniformly for $x$ in a compact subset avoiding the coordinate poles,
\[
 E_n^+(x)=O(n^{-1}),\quad E_n^-(x)=O(n^{-2}),\quad
 E_n^\sigma(x)-E_n^\sigma(0)=O(n^{-3}).
\]
The last estimate follows directly by subtracting the two rational
summands in \eqref{alt:eq:resolvent-tail}. Consequently, if one sign is
negative, the defining series for $\mathcal K$ converges absolutely
normally. In the positive-positive case its difference from the series
at $(0,0)$ is absolutely normally convergent, with summands $O(n^{-3})$.
The scalar series at $(0,0)$ converges by Lemma~\ref{alt:lem:finitecutoff}.
Cauchy's coefficient formula therefore gives
\begin{equation}\label{alt:eq:master-coeff}
 [x^{2a-2}y^{2b-2}]\mathcal K_{\alpha,\beta}
 =-Z(2a,2b-1;\alpha,-\beta)
  -Z(2b,2a-1;\beta,-\alpha)
  -\mathsf L_{2a+2b-1}(c).
\end{equation}

Insert \eqref{alt:eq:parity}. Its two terms
$-\mathsf L_{2a+2b-1}(c)/2$ cancel the diagonal term in
\eqref{alt:eq:master-coeff}. Summation of the remaining binomial terms
uses only the even and odd parts of the binomial theorem. For $N\ge1$,
\begin{align*}
 \sum_{a=1}^{N}\binom{2N}{2a-1}x^{2a-2}y^{2N-2a}
 &=\frac{(x+y)^{2N}-(x-y)^{2N}}{2xy},\\
 \sum_{a=0}^{N}\binom{2N}{2a}x^{2a}y^{2N-2a}
 &=\frac{(x+y)^{2N}+(x-y)^{2N}}2.
\end{align*}
For completeness, the complete expression before replacing $A_\sigma$
by trigonometric functions is
\begin{align}
 \mathcal K_{\alpha,\beta}
 ={}&\left(y^2A_c(y)-\tfrac12\right)U_\alpha
    +\left(x^2A_c(x)-\tfrac12\right)U_\beta\notag\\
 &+\left(A_c(x)-\tfrac1{2x^2}\right)V_{-\beta}
    +\left(A_c(y)-\tfrac1{2y^2}\right)V_{-\alpha}\notag\\
 &+\left(\tfrac1{2x^2}-A_\alpha(x)\right)O_{-\beta}(y)
    +\left(\tfrac1{2y^2}-A_\beta(y)\right)O_{-\alpha}(x).
 \label{alt:eq:master-before-trig}
\end{align}
All the power series used here converge absolutely when
$|x|+|y|<1$. Substitution of $A_\sigma$ proves
\eqref{alt:eq:master}. If the formal constant $T$ is retained in $O_+$,
it cancels directly in this last expression: when $-\beta=1$ one has
$c=\alpha$, and when $-\alpha=1$ one has $c=\beta$.

The left side is holomorphic on
$(\mathbb C\setminus(\mathbb Z\setminus\{0\}))^2$ by the same normal
convergence argument; near a coordinate pole only finitely many tail
terms are singular. The right side is meromorphic on $\mathbb C^2$.
The meromorphic identity theorem extends their equality from the
initial open set. This proves all the removability and continuation
claims in Theorem~\ref{alt:thm:master}. Evenness and exchange symmetry
follow directly from the left side.

There is also a direct check of the coordinate axes. Pair the first
and fifth terms in the braces of \eqref{alt:eq:master} as
\[
 \frac{R_\alpha(x)-R_c(x)}xO_{-\beta}(y)
 +\frac{R_c(x)}x\{O_{-\beta}(y)-V_{-\beta}(x,y)\}.
\]
All $R_\sigma$ have the same principal part $1/(\pi x)$, so the
first quotient is holomorphic at zero. The second difference is
divisible by $x^2$, so the second product is also holomorphic.
The symmetric pair handles $y=0$. Finally, the difference of any
even power evaluated at $x+y$ and $x-y$ is divisible by $xy$,
which proves that $U_\sigma$ is holomorphic on both axes.

\subsection{A diagonal square with elementary rational-parameter values}

\begin{corollary}[A closed alternating square]
\label{alt:cor:square}
For every $x\notin\mathbb Z\setminus\{0\}$,
\begin{equation}\label{alt:eq:diagonal-square}
 \boxed{\displaystyle
 \sum_{n\ge0}(-1)^n(2n+1)
       \left(\sum_{k>n}\frac{(-1)^k}{k^2-x^2}\right)^2
 =-\frac{\pi\{\psi(1+x)+\psi(1-x)+2\gamma\}}
             {4x\sin(\pi x)}.}
\end{equation}
The series is absolutely locally normally convergent. Its value at
$x=0$ is $\zeta(3)/2$.
\end{corollary}
\begin{proof}
Set $\alpha=\beta=-1$ and $y=x$ in
\eqref{alt:eq:master}. Here $c=-1$, $O_+(0)=0$, and
\[
 O_-(2x)-O_-(0)=O_+(x)-O_+(2x).
\]
The two terms at $2x$ cancel, leaving
$\pi O_+(x)/(2x\sin\pi x)$, which is
\eqref{alt:eq:diagonal-square}. At zero,
$O_+(x)=\zeta(3)x^2+O(x^4)$ and
$x\sin(\pi x)=\pi x^2+O(x^4)$. This proves the value and the
removable interpretation.
\end{proof}

The complete diagonal can be written as one symmetric matrix. With
$c_x=\cos(\pi x)$ and the signs ordered as $+,-$,
\begin{equation}\label{alt:eq:diagonal-matrix}
 \bigl(\mathcal K_{\alpha,\beta}(x,x)\bigr)_{\alpha,\beta\in\{+,-\}}
 =\frac{\pi}{2x\sin(\pi x)}
 \begin{pmatrix}
 2c_xO_-(x)-O_+(x)+2\log2 & O_-(x)+c_x\log2\\
 O_-(x)+c_x\log2 & O_+(x)
 \end{pmatrix}.
\end{equation}
This follows from the same substitution $y=x$ and
$O_+(2x)+O_-(2x)=O_+(x)-\log2$ in the other sign cases.
In particular,
\begin{equation}\label{alt:eq:half-matrix}
 \bigl(\mathcal K_{\alpha,\beta}(1/2,1/2)\bigr)_{\alpha,\beta}
 =\pi\begin{pmatrix}1&-1\\-1&2\log2-1\end{pmatrix}.
\end{equation}
Here $E_n^+(1/2)=1/(n+1/2)$, so the positive-positive entry also
reduces directly to the ordinary alternating arctangent series for
$\pi$. This example lies on the apparent oblique polar line $x+y=1$
of the unsimplified master formula, and makes its removability explicit.

For $x=r/N$, $1\le r<N$, Gauss's digamma formula
\cite[Eq.~(5.4.19)]{DLMFGamma} gives the explicit
finite expression
\begin{equation}\label{alt:eq:rational-O}
 O_+(r/N)=\log(2N)-\frac{N}{2r}
 -2\sum_{j=1}^{\lfloor(N-1)/2\rfloor}
       \cos\frac{2\pi rj}{N}\log\sin\frac{\pi j}{N}.
\end{equation}
Thus every such square is $\pi$ times an explicit linear combination,
with algebraic coefficients, of $1$ and logarithms of positive
algebraic numbers. For example,
\begin{align}
 \sum_{n\ge0}(-1)^n(2n+1)
  \left(\sum_{k>n}\frac{(-1)^k}{k^2-1/4}\right)^2
 &=\pi(2\log2-1),\label{alt:eq:half-example}\\
 \sum_{n\ge0}(-1)^n(2n+1)
  \left(\sum_{k>n}\frac{(-1)^k}{k^2-1/9}\right)^2
 &=\frac{3\sqrt3\pi}{2}(\log3-1),\label{alt:eq:third-example}\\
 \sum_{n\ge0}(-1)^n(2n+1)
  \left(\sum_{k>n}\frac{(-1)^k}{k^2-1/16}\right)^2
 &=2\sqrt2\pi(3\log2-2).\label{alt:eq:quarter-example}
\end{align}
The squares are algebraic squares throughout; no complex conjugation
is inserted when $x$ is complex.

Expanding \eqref{alt:eq:diagonal-square} produces a short all-order
convolution. With $\eta(s)=(1-2^{1-s})\zeta(s)$ for $s>1$, for every
$m\ge2$ one has
\begin{equation}\label{alt:eq:convolution}
 \sum_{a=1}^{m-1}\sum_{n\ge0}(-1)^n(2n+1)
        T_{2a}^-(n)T_{2m-2a}^-(n)
 =\frac12\zeta(2m-1)
   +\sum_{j=1}^{m-2}\eta(2j)\zeta(2m-2j-1).
\end{equation}
Indeed
$\pi x/\sin\pi x=1+2\sum_{j\ge1}\eta(2j)x^{2j}$;
coefficient extraction is legitimate by absolute normal convergence.

\begin{corollary}[An anchored Gamma primitive]
\label{alt:cor:gamma-primitive}
For $|u|<1$,
\begin{equation}\label{alt:eq:gamma-primitive}
 \int_0^u x\sin(\pi x)\mathcal K_{-,-}(x,x)\,dx
 =-\frac\pi4\left\{
 \log\Gamma(1+u)-\log\Gamma(1-u)+2\gamma u\right\}.
\end{equation}
The logarithms are their analytic branches on the unit disk,
normalized by $\log\Gamma(1)=0$. The path of integration may be any
path in that disk, and the defining series for $\mathcal K_{-,-}$
may be integrated term by term. In particular the value at $u=1/2$
is $\pi(\log2-\gamma)/4$.
\end{corollary}
\begin{proof}
Multiply \eqref{alt:eq:diagonal-square} by $x\sin(\pi x)$.
Its right side is the derivative of the expression in braces
in \eqref{alt:eq:gamma-primitive}, multiplied by $-\pi/4$.
The anchored logarithms exist because $\Gamma(1\pm u)$ are
holomorphic and nonzero on the simply connected unit disk.
Both sides vanish at zero. Absolute local normal convergence,
already proved, justifies exchanging the path integral and sum.
Finally $\Gamma(3/2)/\Gamma(1/2)=1/2$ gives the stated value.
\end{proof}

Some individual coefficients of the full two-variable theorem are
also useful. Let
$C_{p,q}^{\alpha,\beta}=\sum_{n\ge0}(-1)^n(2n+1)
T_p^\alpha(n)T_q^\beta(n)$. Then
\begin{align}
 C_{2,2}^{+,+}&=\frac{\pi^2\log2}{2}-\frac54\zeta(3),
 & C_{2,2}^{+,-}&=-\frac{\pi^2\log2}{4}-\frac38\zeta(3),
 \label{alt:eq:coeff22}\\
 C_{2,2}^{-,-}&=\frac12\zeta(3),
 &C_{2,4}^{+,+}&=\frac{\pi^4\log2}{48}
                   +\frac{\pi^2\zeta(3)}{12}-\frac{23}{32}\zeta(5),
 \label{alt:eq:coeff24}\\
 C_{4,4}^{+,+}&=-\frac{335}{64}\zeta(7)
                 +\frac{13\pi^2}{48}\zeta(5)
                 +\frac{\pi^4}{32}\zeta(3).
 \label{alt:eq:coeff44}
\end{align}
Only $C_{2,2}^{+,+}$ is conditionally convergent. As an example with a
higher odd centered weight, \eqref{alt:eq:weighted-tail} gives the
absolutely convergent identity
\begin{equation}\label{alt:eq:cubicweight}
 \sum_{n\ge0}(-1)^n\{n^3+(n+1)^3\}T_4^+(n)^2
 =\frac{\pi^4\log2}{24}+\frac{\pi^2\zeta(3)}8
       -\frac{49}{16}\zeta(5).
\end{equation}

\subsection{All repeated rational resolvents}

The same identity supplies explicit derivatives of every order. Define
\[
 E_{n,r}^{\sigma}(x)=\sum_{k>n}
       \frac{\sigma^k}{(k^2-x^2)^{r+1}},\qquad r\ge0,
\]
and let $\partial_{x^2}=(2x)^{-1}\partial_x$ on even analytic
functions, with its removable meaning at $x=0$.

\begin{corollary}[Repeated-resolvent products]\label{alt:cor:derivatives}
For integers $r,s\ge0$ and parameters outside the coordinate poles,
\begin{equation}\label{alt:eq:derivatives}
 \sum_{n\ge0}(-1)^n(2n+1)
        E_{n,r}^{\alpha}(x)E_{n,s}^{\beta}(y)
 =\frac{1}{r!s!}
      \partial_{x^2}^{\,r}\partial_{y^2}^{\,s}
                         \mathcal K_{\alpha,\beta}(x,y).
\end{equation}
The right side is an explicit finite expression obtained by
differentiating \eqref{alt:eq:master}. Whenever $r+s>0$, the series
is absolutely locally normally convergent for all sign pairs.
\end{corollary}
\begin{proof}
Termwise differentiation of the individual rational tails gives
$\partial_{x^2}^{\,r}E_n^\alpha(x)=r!E_{n,r}^\alpha(x)$.
For compact parameter sets avoiding poles, the positive-sign estimate
is $E_{n,r}^+(x)=O(n^{-2r-1})$; an alternating sign improves it to
$O(n^{-2r-2})$. Thus a derivative of positive total order makes the
outer series absolutely summable, even when both signs are positive.
The locally uniform identity in Theorem~\ref{alt:thm:master} and
Cauchy's differentiation formula justify the operation. Evenness
gives the stated removable interpretation at the axes.
\end{proof}

\subsection{A colored polylogarithmic function and its anchored primitive}

The same finite-difference mechanism has an interior-variable form that
also fixes an antiderivative. For $p,q\ge2$, $|a|,|b|\le1$, and
$|z|<1$, put
$T_p^a(n)=\operatorname{Li}_p(a)-\sum_{k=1}^n a^k/k^p$.
Use the colored convention
\[
 \operatorname{Li}_{s_1,\ldots,s_d}(z_1,\ldots,z_d)
 =\sum_{n_1>\cdots>n_d\ge1}
       \frac{z_1^{n_1}\cdots z_d^{n_d}}
            {n_1^{s_1}\cdots n_d^{s_d}}.
\]

\begin{proposition}[Colored generating function and primitive]
\label{alt:prop:GF}
Writing $F(z)=\sum_{n\ge0}T_p^a(n)T_q^b(n)z^n$, one has
\begin{align}
 (1-z)F(z)={}&\operatorname{Li}_p(a)\operatorname{Li}_q(b)
 -\operatorname{Li}_q(b)\operatorname{Li}_p(az)
 -\operatorname{Li}_p(a)\operatorname{Li}_q(bz)\notag\\
 &+\operatorname{Li}_{p,q}(az,b)
  +\operatorname{Li}_{q,p}(bz,a)
  +\operatorname{Li}_{p+q}(abz).
 \label{alt:eq:GF}
\end{align}
The primitive normalized to vanish at zero is
\begin{align}
 \int_0^zF(t)\,dt={}&
  \operatorname{Li}_p(a)\operatorname{Li}_q(b)\operatorname{Li}_1(z)
 -\operatorname{Li}_q(b)\operatorname{Li}_{1,p}(z,a)
 -\operatorname{Li}_p(a)\operatorname{Li}_{1,q}(z,b)\notag\\
 &+\operatorname{Li}_{1,p,q}(z,a,b)
  +\operatorname{Li}_{1,q,p}(z,b,a)
  +\operatorname{Li}_{1,p+q}(z,ab).
 \label{alt:eq:primitive}
\end{align}
The integral is along any path in the unit disk. Every function in the
formulas is defined by an absolutely locally normally convergent series.
\end{proposition}
\begin{proof}
The identity
\[
 T_p^a(n-1)T_q^b(n-1)-T_p^a(n)T_q^b(n)
 =\frac{a^n}{n^p}T_q^b(n)+\frac{b^n}{n^q}T_p^a(n)
   +\frac{(ab)^n}{n^{p+q}}
\]
proves \eqref{alt:eq:GF} after multiplication by $z^n$ and summation.
All rearrangements are absolute in the disk. The general derivative
\[
 \frac{d}{dz}\operatorname{Li}_{1,s_1,\ldots,s_d}
                (z,c_1,\ldots,c_d)
 =\frac{\operatorname{Li}_{s_1,\ldots,s_d}
                (zc_1,c_2,\ldots,c_d)}{1-z}
\]
follows by summing the largest index first. It proves the displayed
primitive term by term. Each term vanishes at $z=0$, which fixes its
additive constant and branch.
\end{proof}
The generating function and primitive are finite-difference consequences,
not assertions of general depth reduction. The additional reduction in
Theorem~\ref{alt:thm:master} uses the specific alternating weight and
odd-weight parity input. It does not imply the outstanding Gaussian
$S_6$ or revised $S_8$ identities.

\section{Source reconciliation and a further research programme}
\label{sec:audit}

\subsection{What the audit establishes}

The source claims used in the proofs were checked against the pinned
revision, including the finite-cutoff conventions and the exact matrices.
The audit establishes one stale research-status statement and a precise
centering reconciliation. It did not establish an additional false theorem
in the inspected parts of the six incoming reports. This conclusion is
limited to those parts: acquiring the full source inventory does not
constitute checking every proof in the collective manuscript.

\begin{longtable}{@{}>{\raggedright\arraybackslash}p{0.26\textwidth}>{\raggedright\arraybackslash}p{0.38\textwidth}>{\raggedright\arraybackslash}p{0.28\textwidth}@{}}
\toprule
Source & Material checked for this continuation & Result \\
\midrule
\endfirsthead
\toprule
Source & Material checked for this continuation & Result \\
\midrule
\endhead
Canonical fifth-moment report \cite{Canonical} &
Tail blocks, inverse transform, diagonal rank, and the first-$q$ conjecture &
Normalization retained; gap four advances its conjecture \\
Differential polygamma and Lerch jets \cite{Differential} &
Adjacent-order proof and fixed-gap recurrence &
Recurrence and determinant signs retained \\
Bernoulli--Tornheim--digamma \cite{Bernoulli} &
Fixed smaller orders $2,4,6$; shifted-Hankel reduction; tail foundation;
Tornheim subtraction; differential and research sections &
No additional false formula found; R8 status needs reconciliation \\
Polygamma tail identities \cite{Polygamma} &
Moment rank and reconstruction; spectral polynomial formulas; sixth
mixed moment &
Supplies the already available spectral coordinates used below \\
Cubic tail descent \cite{Cubic} &
Cubic Gamma product and the master identity with its subtraction &
No error established in the inspected product and limit arguments \\
Determinantal harmonic resonance \cite{Resonance} &
First resonance and the cutoff-to-spectral transfer law &
The amplitude-derivative correction remains necessary \\
Multishift reflection \cite{Multishift} &
Entire-order transform and full horizontal-cut reflection &
No error established in the inspected holomorphy and jump conventions \\
\bottomrule
\end{longtable}

The accompanying source-audit record identifies the exact section files
and equation labels. The preserved Canonical report is already under
\path{Analysis/Polylogarithms/docs/reports/stieltjes-harmonic-resonance/continuations/canonical-digamma-fifth/};
it is not one of the six current incoming archives. The earlier extracted
tail-rank source used during this work was checked against that preserved
file and has the identical Git blob.

\subsection{A sixth mixed moment whose spectral evaluation already exists}
\label{sh:audit-sixth}

Research item R8 of the Bernoulli--Tornheim--digamma report asks for a
harmonic-spectral evaluation of
\[
 \widetilde M_6=\FP\int_0^1(\psi(x)+\gamma)^2\psi'(x)^2\,dx.
\]
Here and throughout this subsection, $\FP\int_0^1f(x)\,dx$ means the
coefficient of $\varepsilon^0(\log\varepsilon)^0$ in
$\int_\varepsilon^1f(x)\,dx$ as $\varepsilon\downarrow0$, in the
original coordinate $x$. The operation is linear on the Laurent-log
expansions used here. The existence question is already answered by
the Polygamma Tail Identities report \cite{Polygamma}. Its equation
\texttt{sixth:I} evaluates the uncentered mixed integral; equations
\texttt{spectral:exampleDP} and
\texttt{spectral:digamma\_trigammasquare} evaluate the lower mixed term.
The differing centering can be reconciled explicitly.

Define, in decreasing-index convention,
\begin{align}\label{audit:DP}
 D_P(s)={}&-4\zeta(3,s+1)-4\zeta(s+4)
          +2\zeta(s+2,2)+2\zeta(2)\zeta(s+2)\notag\\
         &+4\zeta(s+3,1)-4\gamma\zeta(s+3),\qquad \Re s>-1.
\end{align}
All the displayed series are holomorphic in this half-plane. To specify
the other inherited coordinate without leaving an unnamed integral,
put, for $n\ge0$,
\begin{align*}
 b_n&=H_n-\gamma,&c_n&=H_n^{(2)}+\zeta(2),&
 d_n&=H_n^{(3)}-\zeta(3),\\
 e_n&=H_n^{(4)}+\zeta(4),&f_n&=H_n^{(5)}-\zeta(5),
\end{align*}
and
\begin{align}
 r_n={}&4b_n^2d_n+2b_nc_n^2-6c_nd_n-10b_ne_n+6f_n,
       \label{audit:residue}\\
 \Omega_5={}&-\sum_{n\ge1}(r_n-r_{n-1})\log n
           =\sum_{n\ge1}r_n\log(1+1/n).\label{audit:Omega}
\end{align}
These are ordinary absolutely convergent series. For example, the
harmonic-tail expansions and $8\zeta(2)^2=20\zeta(4)$ give
$r_n=O(\log(n+1)/n+\log^2(n+1)/n^2)$ and
$r_n-r_{n-1}=O(\log(n+1)/n^2)$; summation by parts has a vanishing
boundary term. These definitions agree exactly with the source's
\texttt{sixth:D} and \texttt{sixth:omega-series}.

\begin{proposition}[Reconciliation of the centered sixth coordinate]
\label{audit:sixth}
The inherited spectral evaluations imply
\begin{equation}\label{sh:centered-sixth-reconciliation}
 \boxed{\widetilde M_6
   =\Omega_5+2\gamma D_P'(0)+4\gamma^2\zeta'(3)
      +\frac{31}{3}\zeta(2)\zeta(3)-\frac{19}{3}\zeta(5).}
\end{equation}
\end{proposition}
\begin{proof}
The three source evaluations in the same finite-part convention are
\begin{align*}
 \FP\int_0^1\psi^2\psi'^2\,dx
 &=\Omega_5+2\gamma^2\zeta(3)-\frac{35}{3}\gamma\zeta(4)
       +\frac{31}{3}\zeta(2)\zeta(3)-\frac{19}{3}\zeta(5),\\
 \FP\int_0^1\psi\psi'^2\,dx
 &=\frac{35}{6}\zeta(4)-2\gamma\zeta(3)+D_P'(0),\\
 \FP\int_0^1\psi'^2\,dx
 &=2\zeta(3)+4\zeta'(3).
\end{align*}
Expand $(\psi+\gamma)^2\psi'^2$ and use linearity. The coefficients
of $\gamma\zeta(4)$ and $\gamma^2\zeta(3)$ cancel exactly, leaving
\eqref{sh:centered-sixth-reconciliation}. In the Canonical fifth-moment
notation one also has $D_P(s)=2J(s+1)$, by the stuffle product; hence
$D_P'(0)=2J'(1)$.
\end{proof}

The proposed correction to R8 is therefore: determine whether the
specified combination in \eqref{sh:centered-sixth-reconciliation}
reduces to a smaller, explicitly stated collection of constants, such
as selected ordinary zeta derivatives and prescribed height-one
spectral data. Existence of a harmonic-spectral representation should
be marked resolved. This reconciliation does not prove a reduction
to that smaller collection, and it is not an arithmetic-independence
statement about $\Omega_5$.

\subsection{Conventions that must survive integration}

Three distinctions materially affect the values of the formulas.
First, the original centered lattice has
$\zeta(-2j,1/2)=0$, whereas a shifted lattice generally does not.
Theorems~\ref{sh:transform} and \ref{sh:jets} give the two required
corrections explicitly. Second, equality of zero-order continued and
ordinary values does not imply equality after spectral differentiation:
the nonzero $\zeta^{(j)}(-k,\alpha)$ terms must be retained. Third,
$\mathcal K_{+,+}$ and its $(2,2)$ coefficient have an ordinary
conditional interpretation; their original series must not be
described as absolutely convergent. The other sign pairings are
absolutely convergent.

These are normalization requirements for using the new results, not
claims that the inspected predecessors made all three mistakes.
The separately unresolved Gaussian $S_6$ and revised $S_8$ candidates
remain conjectural in this continuation. A proved reduction for a
different tail family supplies no missing identity for those values.

\subsection{Further research questions with concrete targets}
\label{sec:questions}

The following questions concern exact evaluations, structure, or proof
certificates. The numerical evidence accompanying the sources can help
choose cases, but each proposed endpoint is an identity or a rigorous
nonvanishing theorem.

\begin{enumerate}[label=\textbf{Q\arabic*.},ref=Q\arabic*,itemsep=10pt]
\item \textbf{Complete the optimal-prefix conjecture beyond gap four.}
For even $8\le p<q$ with $q-p\ge6$, prove or refute
$\det M_{p,q}\ne0$. The first unproved fixed-gap family is
$(p,q)=(2b-6,2b)$ with $b\ge7$. Its boundary determinant has size
three. The explicit sign change already seen between $b=6$ and $b=7$
rules out a proof based on one fixed sign for this family. A useful
target is an exact expression whose factors prove nonvanishing while
allowing the determinant's sign to vary. Comparing the size-$c$
fixed-gap determinant with the size-$p/2$ complementary determinant
in the shifted-Hankel approach may reduce different parameter ranges.

\item \textbf{Classify higher minors of the fundamental recurrence.}
Theorem~\ref{eg:two-minors} is uniform in the gap and the two chosen
columns. Determine which larger minors vanish or change sign, and
whether a recurrence for their exterior products separates those two
possibilities. The one-negative-index induction in the present proof
does not directly extend to several sign changes. A successful
replacement should explain the recorded third-order sign reversal,
not impose a positivity property contradicted by it.

\item \textbf{Determine the exact formal rank after a Hurwitz shift.}
Section~\ref{sh:section} gives at most $q+2$ coordinates over
$\Q(v)$. Is that list generically minimal, and how many initial
even moments reconstruct it? Describe every exceptional shift in a
stated real or complex domain, including the known degeneration at
$v=1$. This is a question about rational coefficient matrices; it
does not require independence of the specialized Hurwitz values.
The Bernoulli-polynomial endpoint and the highest single-Hurwitz
column should be kept throughout any determinant computation.

\item \textbf{Allow different shifts in the two factors.}
Replace the product by
$\zeta(p,n+v)\zeta(q,n+w)$. Integer values of $v-w$ can be
related by finite tail corrections. For general rational differences,
determine an explicit finite normal form and the necessary depth-two
shifted coordinates. The first test is whether a discrete primitive
analogous to Theorem~\ref{sh:primitive} closes in single Hurwitz
values and endpoint products, and exactly which shift differences
permit that closure.

\item \textbf{Combine tail-order derivatives with the weight jets.}
Theorem~\ref{sh:jets} differentiates the factor $x_n^{-s}$.
Derivatives in $p$ or $q$ instead create logarithmic harmonic tails.
Construct the mixed Taylor-Laurent expansion in all three variables,
including motion of the polar hyperplanes and coincidences of
resonances. The target is a finite expression in specified derivatives
of single and multiple Hurwitz functions and generalized Stieltjes
constants, with the scale and amplitude corrections visible at every
order. This would extend the source's Lerch and Tornheim jet machinery
to the shifted product setting.

\item \textbf{Replace signs by higher roots of unity.}
For a fixed conductor $N$, study the analogues of
$\mathcal K_{\alpha,\beta}$ with $N$th roots of unity as colors
and a corresponding outer character. State explicitly which
conjugation and weight-parity sectors reduce in depth. The level-two
parity formula used here cannot simply be reused unchanged. A
concrete objective is a conductor-four generator expressed through
digamma at quarter shifts and a specified collection of cyclotomic
polylogarithms. Such a result would give an exact bridge toward the
Gaussian setting without assuming the outstanding Gaussian identities.

\item \textbf{Find compact normal forms for differentiated resolvents.}
Parameter differentiation of Theorem~\ref{alt:thm:master} already
gives sums of products of arbitrary repeated rational resolvents.
Can their right sides be organized by finite Bell-polynomial formulas
that display the removable singularities before specialization?
At rational parameters, seek explicit formulas in algebraic
coefficients, logarithms, and Dirichlet $L$-values arising from
polygamma at rational arguments. A separate useful target is a direct
contour or hypergeometric proof of the diagonal square and its Gamma
primitive, independent of the general parity theorem.

\item \textbf{Sum the higher odd weights into one generating identity.}
The finite reductions in \eqref{alt:eq:weighted-tail} handle every
admissible odd polynomial weight. Introduce a third variable for
that degree and determine a compact generating function, or an
explicit differential operator on $O_\pm$. The admissible degree
condition must be built into the formulation: extending it crosses
new convergence thresholds and requires additional subtractions.
An exact endpoint prescription is part of the desired identity.

\item \textbf{Extend the discrete primitive to higher products.}
For three or more Hurwitz tails, finite differences introduce products
of fewer tails. Determine which polynomial weights allow this
hierarchy to terminate in single Hurwitz values, and which force
genuine multiple-Hurwitz coordinates. Hoffman's higher-product
decomposition \cite{Hoffman} supplies the ordered-sum framework.
A useful initial problem is a complete classification for three
equal even orders and weights of degree below the first resonance.

\item \textbf{Reduce the reconciled sixth spectral coordinate.}
Choose a target collection first, then test whether
$\Omega_5+2\gamma D_P'(0)+4\gamma^2\zeta'(3)$ belongs to its
rational span. Candidate collections can include specified ordinary
zeta derivatives and the height-one slopes already defined by the
Gamma bridge. The source gives convergent series and a finite
multiple-zeta derivative representation for the current coordinate;
the missing step is a proved functional reduction among those
representations. Numeric relation searches alone cannot decide it.

\item \textbf{Connect the colored primitive to endpoint periods.}
Proposition~\ref{alt:prop:GF} fixes both an analytic function and its
additive constant in the disk. Continue it to selected boundary
points using a stated path, compute its endpoint limits, and compare
the resulting depth-three polylogarithms under shuffle and stuffle
relations. Determine which combinations admit a Gamma or polygamma
description like the diagonal primitive. At a branch point, the
desired result should include the full finite-part or limiting rule.

\item \textbf{Produce exact certificates for the remaining numerical candidates.}
For the Gaussian $S_6$ and revised $S_8$ conjectures, isolate a finite
algebraic, contour, or hypergeometric identity whose specialization
would prove the claimed value. Separate the convergence argument
from the period reduction, and state the intermediate functions with
fixed branches. The finite telescoper in \eqref{eg:V-difference}
illustrates the desired form of a checkable certificate for a sum;
the Gaussian questions need their own certificates.
\end{enumerate}

The immediate priorities are Q1, Q3, and Q6. They extend the three
proved mechanisms of this article while preserving an exact target
and a clear test for what would constitute a solution.

\section{Reproducibility, proof status, and integration}
\label{sec:reproducibility}

\subsection{Proofs and computational checks have different roles}

The infinite-family conclusions are proved in the article. The sign
comparison in Theorem~\ref{eg:two-minors}, the finite Bernoulli
primitive in Theorem~\ref{sh:primitive}, and the normally convergent
coefficient argument in Theorem~\ref{alt:thm:master} are the respective
all-order arguments. The programs reproduce exact finite instances,
check normalizations by different representations, and provide
numerical diagnostics for analytic specializations. No integer-relation
search is used in any of the supplied verifiers.

\begin{center}
\begin{tabular}{@{}>{\raggedright\arraybackslash}p{0.34\textwidth}>{\raggedright\arraybackslash}p{0.18\textwidth}>{\raggedright\arraybackslash}p{0.4\textwidth}@{}}
\toprule
Verification & Scope & What is compared \\
\midrule
Gap-four determinants & $3\le b\le18$ &
Original Bernoulli matrix, recurrence determinant, and product formula \\
Generated minors & 1,946 minors &
Exact signs for gaps $2c$, $c=2,\ldots,6$ \\
Finite telescopers & $(4,8),(8,12)$ &
Rational-function differences and the complete constant-term limit \\
Shifted discrete primitives & 168 triples &
Generic-variable identities, $2\le p\le q\le8$, $q\le d\le q+5$ \\
Shifted transforms & $21+21$ checks &
Inverse Bernoulli polynomials and independent finite rational endpoints \\
Shift derivatives & 23 coefficients &
The closed convolution for derivative orders $2,\ldots,24$ \\
Alternating master & 144 coefficients &
Every even $p,q\ge2$, $p+q\le18$, for all four sign pairs \\
Formal endpoint cancellation & 256 checks &
Cancellation of the single formal $\zeta(1)$ symbol \\
Analytic diagnostics & $36+16$ checks &
Independent shifted asymptotics and signed digamma-tail sums \\
\bottomrule
\end{tabular}
\end{center}

The gap-four program also verifies the row and column scalings
against direct cotangent coefficient extraction, checks the highest
polynomial coefficient, and reproduces the higher-gap sign obstruction.
The alternating program expands the meromorphic right side as a
Laurent polynomial in two variables, verifies cancellation of axial
negative powers and odd powers, and checks the diagonal convolution.
The shifted program treats the relevant Hurwitz values as formal
symbols obeying their exact recurrence, so its discrete-primitive
checks are symbolic identities in the shift, not tests at one shift.

The shifted diagnostics use 85 decimal digits and the alternating
diagnostics use 55. Each comparison has a $10^{-45}$ acceptance
threshold. The observed residuals are below $7.1\times10^{-65}$ and
$6.6\times10^{-55}$, respectively, in the recorded run. These numbers
are numerical consistency measurements, not certified error
enclosures. The proofs, including their convergence arguments, do not
depend on treating those measurements as exact zeros.

\subsection{Contents of the accompanying archive}

The archive includes a standalone TeX source and its compiled PDF,
the modular source used for editing, all three independent verification
programs, their complete JSON results, and a reproducibility record.
The source manifest fixes the repository revision, paths, byte sizes,
Git blob identifiers, and local SHA-256 checksums for the acquired
inputs. The proof-status record distinguishes new theorems,
inherited ingredients, finite computational checks, and unresolved
questions. No large predecessor PDF or redundant copy of the
repository is required to rebuild this article.

For verification from the archive directory, the three commands are
\begin{quote}\small\ttfamily
python3 verification/verify\_gap\_four.py\\
\hspace*{1em}--output results/gap\_four\_exact\_results.json\\[4pt]
python3 verification/verify\_shifted\_tails.py\\
\hspace*{1em}--output results/shifted\_tail\_results.json\\[4pt]
python3 verification/verify\_alternating\_master.py --numeric\\
\hspace*{1em}--output results/alternating\_verification.json
\end{quote}
Each indented continuation belongs to the preceding shell command.
The plain-text README contains directly copyable single-line commands.
Python, SymPy, and mpmath are sufficient for these checks. The exact
dependency versions used are recorded in the package. Rebuilding the
document requires a LaTeX installation with the packages declared in
its preamble; the provided build program assembles the standalone
source and runs \texttt{latexmk}.

\subsection{Suggested integration points}

The gap-four theorem belongs beside the optimal-prefix conjecture and
the existing adjacent-order and fixed-smaller-order proofs. Preserve
the matrix convention in \eqref{eg:entries}, especially the omission
of a binomial coefficient with lower index zero. Record the newly
settled family as $q-p=4$ and leave Q1's remaining range open.

The Hurwitz-shift transform can follow the centered-tail foundation.
Keep \eqref{sh:Dhat} and \eqref{sh:ordinarycorrection} close to the
transform itself, because they are part of its value normalization.
The all-jet theorem belongs with the spectral-subtraction and Lerch
jet material; its derivatives act on the weight exponent.

The signed master, Gamma primitive, and colored polylogarithmic
primitive form a self-contained continuation of the tail generating
functions. Retain the explicit convergence distinction for the
positive-positive entry. Finally, update the sixth mixed-moment
research item using Proposition~\ref{audit:sixth}; retain the
smaller-collection reduction problem as Q10.

These integration suggestions are supplied as reviewable material.
The repository itself was not modified during this research.

\begin{thebibliography}{99}
\raggedright
\addcontentsline{toc}{section}{References}

\bibitem{Repo}
Vladimir Reshetnikov, maintainer,
\emph{ProveIt}, repository snapshot
\texttt{446dca9bd65c20e409d1d2561d4ffad408103e2a},
11 October 2026.
\href{https://github.com/VladimirReshetnikov/ProveIt/tree/446dca9bd65c20e409d1d2561d4ffad408103e2a/Analysis/Polylogarithms/docs/manuscript}{Polylogarithm manuscript sources};
\href{https://github.com/VladimirReshetnikov/ProveIt/tree/446dca9bd65c20e409d1d2561d4ffad408103e2a/docs/incoming}{six incoming reports}.
The accompanying manifest identifies every acquired source blob.

\bibitem{Canonical}
ProveIt research reports,
\emph{Canonical Digamma and Fifth Moment},
preserved \texttt{canonical-digamma-fifth} continuation,
11 October 2026. In particular,
\texttt{sections/07\_tail\_rank.tex} and its optimal-prefix conjecture.
\href{https://github.com/VladimirReshetnikov/ProveIt/tree/446dca9bd65c20e409d1d2561d4ffad408103e2a/Analysis/Polylogarithms/docs/reports/stieltjes-harmonic-resonance/continuations/canonical-digamma-fifth}{Preserved report at the pinned revision}.

\bibitem{Differential}
ProveIt research reports,
\emph{Differential Polygamma and Lerch Identities},
11 October 2026. In particular,
\texttt{sections/08\_adjacent\_tails.tex}.
\href{https://github.com/VladimirReshetnikov/ProveIt/blob/446dca9bd65c20e409d1d2561d4ffad408103e2a/docs/incoming/ProveIt_Differential_Polygamma_and_Lerch_Jets_2026-10-11.zip}{Incoming source archive}.

\bibitem{Bernoulli}
ProveIt research reports,
\emph{Bernoulli Determinants, Cubic Tornheim Jets, and Canonical
Digamma Integrals}, 11 October 2026.
The tail foundation, general shifted-Hankel reduction,
fixed smaller orders four and six, and research item R8.
\href{https://github.com/VladimirReshetnikov/ProveIt/blob/446dca9bd65c20e409d1d2561d4ffad408103e2a/docs/incoming/ProveIt_Bernoulli_Tornheim_Digamma_2026-10-11.zip}{Incoming source archive}.

\bibitem{Polygamma}
ProveIt research reports,
\emph{Exact Polylogarithm Parity, Harmonic-Tail Reconstruction,
and Higher Polygamma Moments}, 11 October 2026.
In particular, \texttt{sections/07\_spectral\_polynomials.tex}
and \texttt{sections/08\_sixth\_moment.tex}.
\href{https://github.com/VladimirReshetnikov/ProveIt/blob/446dca9bd65c20e409d1d2561d4ffad408103e2a/docs/incoming/ProveIt_Polygamma_Tail_Identities_2026-10-11.zip}{Incoming source archive}.

\bibitem{Cubic}
ProveIt research reports,
\emph{Cubic Invariant Tail Descent}, 10 October 2026.
Cubic Gamma product and master identity with explicit subtraction.
\href{https://github.com/VladimirReshetnikov/ProveIt/blob/446dca9bd65c20e409d1d2561d4ffad408103e2a/docs/incoming/ProveIt_Cubic_Tail_Descent_2026-10-10.zip}{Incoming source archive}.

\bibitem{Resonance}
ProveIt research reports,
\emph{Determinantal Harmonic Resonance}, 10 October 2026.
First resonance and finite-part transfer; archive root
\texttt{ProveIt\_Determinantal\_Harmonic\_Resonance/}.
\href{https://github.com/VladimirReshetnikov/ProveIt/blob/446dca9bd65c20e409d1d2561d4ffad408103e2a/docs/incoming/ProveIt_Determinantal_Harmonic_Resonance_2026-10-10.zip}{Incoming source archive}.

\bibitem{Multishift}
ProveIt research reports,
\emph{Contact-Complete Multishift Lerch Reflection},
10 October 2026. Entire-order transform and horizontal cuts.
\href{https://github.com/VladimirReshetnikov/ProveIt/blob/446dca9bd65c20e409d1d2561d4ffad408103e2a/docs/incoming/ProveIt_Multishift_Reflection_2026-10-10.zip}{Incoming source archive}.

\bibitem{Hoffman}
M.~E.~Hoffman,
\emph{Sums of Products of Riemann Zeta Tails},
Mediterranean Journal of Mathematics \textbf{13} (2016), 2771--2781.
\url{https://arxiv.org/abs/1602.03399}.
\url{https://doi.org/10.1007/s00009-015-0653-9}.

\bibitem{Teo}
L.-P.~Teo,
\emph{Alternating Double Euler Sums, Hypergeometric Identities and a
Theorem of Zagier},
Journal of Mathematical Analysis and Applications \textbf{462}
(2018), 777--800.
\url{https://arxiv.org/abs/1705.01269}.
\url{https://doi.org/10.1016/j.jmaa.2018.02.037}.

\bibitem{XuWang}
C.~Xu and W.~Wang,
\emph{Explicit formulas of Euler sums via multiple zeta values},
arXiv:1805.08056v3, 2019, equation (3.7).
\url{https://arxiv.org/abs/1805.08056}.

\bibitem{DilcherJiu}
K.~Dilcher and L.~Jiu,
\emph{Hankel Determinants of Sequences Related to Bernoulli and Euler
Polynomials}, arXiv:2007.09821, 2020.
\url{https://arxiv.org/abs/2007.09821}.

\bibitem{DLMF}
NIST Digital Library of Mathematical Functions,
\emph{Hurwitz Zeta Function}, Section 25.11.
Recurrence, polygamma relation, negative integer values,
and Euler--Maclaurin representations.
\url{https://dlmf.nist.gov/25.11}.

\bibitem{DLMFBernoulli}
NIST Digital Library of Mathematical Functions,
\emph{Bernoulli and Euler Polynomials: Basic Properties}, Section 24.4.
\url{https://dlmf.nist.gov/24.4}.

\bibitem{DLMFGamma}
NIST Digital Library of Mathematical Functions,
Sections 4.22, 5.4, 5.7, and 5.15.
Cotangent partial fractions, Gauss's digamma formula, Taylor series,
and polygamma identities.
\url{https://dlmf.nist.gov/4.22};
\url{https://dlmf.nist.gov/5.4};
\url{https://dlmf.nist.gov/5.7};
\url{https://dlmf.nist.gov/5.15}.

\end{thebibliography}

\end{document}
